% BEGIN EMBEDDED BIBLIOGRAPHY
% Bibliography data are part of this source; the compiler writes its temporary input.
\begin{filecontents*}[overwrite]{\jobname.bib}
@unpublished{QED3Manuscript,
  title = {The field-strength {Osterwalder--Schrader} axioms for {QED} in three dimensions},
  date = {2026},
  shorthand = {Cur},
  note = {Comprehensive edition, 263 pp., 20 October 2026, before the modular integration revision}
}
@unpublished{MaxwellBlockManuscript,
  title = {Maxwell block minimization and {Gaussian} fluctuations},
  date = {2026},
  shorthand = {Max},
  note = {Companion manuscript}
}
@online{EtingofLieNotes,
  author = {Etingof, Pavel},
  title = {Lie Groups and Lie Algebras},
  organization = {Massachusetts Institute of Technology},
  date = {2024},
  note = {Course notes, Section 22.3, Theorem 22.14},
  url = {https://ocw.mit.edu/courses/18-755-lie-groups-and-lie-algebras-ii-spring-2024/mit18_755_s24_lec_full.pdf}
}
\end{filecontents*}
% END EMBEDDED BIBLIOGRAPHY
% FBM1: full proof companion to the separately maintained FBO1 outline.
\documentclass[11pt,letterpaper,reqno]{amsart}
% BEGIN INLINE manuscript-style
\makeatletter
\NeedsTeXFormat{LaTeX2e}[2020/10/01]
% Use with \documentclass[11pt,letterpaper,reqno]{amsart}.
% This file controls presentation only. It supplies no mathematical content.

\RequirePackage[T1]{fontenc}
\RequirePackage{lmodern}
\RequirePackage{microtype}
\RequirePackage[letterpaper,margin=1in,headheight=15pt]{geometry}
\RequirePackage{amsmath,amsthm,amssymb,mathtools}
\RequirePackage{enumitem}
\RequirePackage{perpage}
\RequirePackage{xcolor}
\definecolor{manuscriptred}{rgb}{0.6,0.1,0.1}
\definecolor{manuscriptblue}{rgb}{0.1,0.3,0.5}
\definecolor{manuscriptmagenta}{HTML}{60445C}
\RequirePackage[breaklinks=true,colorlinks=true,
  linkcolor=manuscriptred,citecolor=manuscriptblue,urlcolor=manuscriptmagenta]{hyperref}
\RequirePackage{cleveref}
\hypersetup{pdfauthor={}}

% Retain the AMS paragraph, title, heading and running-head conventions.
% Equations retain one global Arabic sequence, as in the reference.
\setcounter{tocdepth}{1}
\renewcommand{\thefigure}{\Roman{figure}}
\renewcommand{\thetable}{\Roman{table}}
\MakePerPage{footnote}
\renewcommand{\thefootnote}{%
  \ifnum\value{footnote}>9\arabic{footnote}\else\fnsymbol{footnote}\fi}

% Upright statements. Main results are A, B, C; supporting statements
% share a section-local sequence. These end marks indicate typography,
% never proof completeness, validation, or a deferred-proof status.
\theoremstyle{definition}
\newtheorem{ms@theorem}{Theorem}
\renewcommand{\thems@theorem}{\Alph{ms@theorem}}
\newtheorem*{ms@theoremstar}{Theorem}
\newtheorem{ms@proposition}{Proposition}[section]
\newtheorem{ms@lemma}[ms@proposition]{Lemma}
\newtheorem{ms@definition}[ms@proposition]{Definition}
\newtheorem{ms@corollary}[ms@proposition]{Corollary}
\newtheorem*{ms@corollarystar}{Corollary}

\newcommand{\ms@majorstart}{%
  \pushQED{\qed}\renewcommand{\qedsymbol}{$\blacksquare\square$}}
\newcommand{\ms@minorstart}{%
  \pushQED{\qed}\renewcommand{\qedsymbol}{$\blacksquare$}}
\newenvironment{theorem}[1][]{\ms@majorstart\ms@theorem[#1]}
  {\popQED\endms@theorem}
\newenvironment{theorem*}[1][]{\ms@majorstart\ms@theoremstar[#1]}
  {\popQED\endms@theoremstar}
\newenvironment{proposition}[1][]{\ms@minorstart\ms@proposition[#1]}
  {\popQED\endms@proposition}
\newenvironment{lemma}[1][]{\ms@minorstart\ms@lemma[#1]}
  {\popQED\endms@lemma}
\newenvironment{definition}[1][]{\ms@minorstart\ms@definition[#1]}
  {\popQED\endms@definition}
\newenvironment{corollary}[1][]{\ms@majorstart\ms@corollary[#1]}
  {\popQED\endms@corollary}
\newenvironment{corollary*}[1][]{\ms@majorstart\ms@corollarystar[#1]}
  {\popQED\endms@corollarystar}
% The source's lemma/proposition variants are deliberately not carried over:
% choose one marker for each statement type rather than encode proof status.

\theoremstyle{remark}
\newtheorem{remark}[ms@proposition]{Remark}
\newtheorem*{remark*}{Remark}
% The ordinary amsthm proof environment retains its open square.

\crefname{ms@theorem}{theorem}{theorems}
\Crefname{ms@theorem}{Theorem}{Theorems}
\crefname{ms@proposition}{proposition}{propositions}
\Crefname{ms@proposition}{Proposition}{Propositions}
\crefname{ms@lemma}{lemma}{lemmas}
\Crefname{ms@lemma}{Lemma}{Lemmas}
\crefname{ms@definition}{definition}{definitions}
\Crefname{ms@definition}{Definition}{Definitions}
\crefname{ms@corollary}{corollary}{corollaries}
\Crefname{ms@corollary}{Corollary}{Corollaries}
\crefname{remark}{remark}{remarks}
\Crefname{remark}{Remark}{Remarks}

% Bibliography configuration for this manuscript.
  \RequirePackage[backend=biber,style=alphabetic,backref=true,
    giveninits=true,isbn=false,date=year]{biblatex}
  \addbibresource{\jobname.bib}
  \setcounter{biburlnumpenalty}{100}
  \setcounter{biburllcpenalty}{100}
  \setcounter{biburlucpenalty}{100}
  \renewcommand*{\bibfont}{\footnotesize}
  \renewbibmacro{in:}{}
  \DeclareFieldFormat{pages}{#1}
  % Use the public alphabetic style. No private author-hash overrides;
  % repeated author names are printed in full.
\makeatother
% END INLINE manuscript-style
\usepackage{mathrsfs}
% Paper-specific notation.
\newcommand{\R}{\mathbb R}
\newcommand{\C}{\mathbb C}
\newcommand{\Z}{\mathbb Z}
\newcommand{\N}{\mathbb N}
\newcommand{\dd}{\mathrm d}
\DeclareMathOperator{\supp}{supp}
\DeclareMathOperator{\tr}{tr}
\title[Chamber fermion locality and local integration]
 {Chamber fermion locality and local integration}
\date{}
\hypersetup{pdftitle={Chamber fermion locality and local integration},pdfauthor={}}
\begin{document}
\begin{abstract}
We prove uniform locality estimates for penalized chamber Dirac
operators on three-dimensional periodic lattices. A gap on the
mean-zero chamber components and a free Schur-complement estimate
give local inverse bounds independent of the mesh and reference
mass. Local real gauges and convergent walk expansions extend these
bounds to complex connections. We obtain analytic coarse covariances,
an exact weighted Grassmann block identity, and a local expansion of
the normalized determinant. Under additional charge and volume
conditions, the local determinant difference gains a factor
\(e_k^{1/4}\).
\end{abstract}
\maketitle
\tableofcontents
\section{Introduction}\label{sec:introduction}
Fix \(0<\rho<1\), an odd integer \(L\ge5\), \(0<\alpha<1\), and
\(\varepsilon>0\). At level \(k\ge0\), put \(\eta=\eta_k=L^{-k}\).
A photon site is \(c\in\eta\Z^3/\ell\Z^3\); a chamber site is
\(x=(c,g)\), where \(g\in O_{\mathrm{cub}}\) is a signed permutation
matrix. There are \(48\) labels \(g\) and four spin components at
each photon site. Connections are periodic complex oriented bond
functions, with \(A(c',c)=-A(c,c')\); spinors are antiperiodic in
each coordinate. The chamber measure is \(\eta^3/48\), while the
photon-site spinor measure is \(\eta^3\).
Operator kernels use the chamber measure in their column variable;
the identity kernel is \(48\eta^{-3}\mathbf1[x=x']1_4\).
The cluster injection and mean are
\[
 (J\xi)(c,g)=\xi(c),\qquad J^*f(c)=48^{-1}\sum_g f(c,g).
\]
Adjoints use these Hermitian pairings. Transposes use the corresponding
bilinear pairings, without complex conjugation.

Use the Pauli and Dirac matrices
\[
 \sigma_1=\begin{pmatrix}0&1\\1&0\end{pmatrix},\quad
 \sigma_2=\begin{pmatrix}0&-i\\i&0\end{pmatrix},\quad
 \sigma_3=\begin{pmatrix}1&0\\0&-1\end{pmatrix},\qquad
 \gamma_j=\sigma_3\otimes\sigma_j,\quad
 \gamma_5=\sigma_1\otimes1_2.
\]
Let \(s_0\) reverse the third coordinate, \(s_1\) interchange
coordinates \(2,3\), and \(s_2\) interchange \(1,2\); let \(r_{e_1}\)
reverse the first coordinate. The four neighbors, inward normals,
and weights are
\[
 \begin{gathered}
 x^{(i)}=(c,gs_i)\quad(0\le i\le2),\qquad
 x^{(3)}=(c+\eta ge_1,gr_{e_1}),\\
 (n_0,n_1,n_2,n_3)
   =(e_3,(e_2-e_3)/\sqrt2,(e_1-e_2)/\sqrt2,-e_1),\\
 (w_0,w_1,w_2,w_3)=(6,6\sqrt2,6\sqrt2,6).
 \end{gathered}
\]
For a charge \(e_k>0\), the chamber operator is
\begin{equation}\label{eq:intro-chamber-operator}
 (D_Af)(x)=-\eta^{-1}\sum_{i=0}^3
       \frac{w_i}{2}(\rho+\gamma\cdot gn_i)
       \bigl(U_A(x,x^{(i)})f(x^{(i)})-f(x)\bigr),
\end{equation}
where internal transports are one and the external transport is
\(\exp(ie_k\eta A(c,c+\eta ge_1))\).

The unit blocks are
\(\Delta_z=\{c:|c-z|_\infty<1/2\}\), \(z\in\Z^3/\ell\Z^3\).
For \(c\in\Delta_z\), let \(c=c_0,c_1,\ldots,c_k=z\) be the
centers of its nested blocks of sides \(L^j\eta\).
Write \(\tau A(u,v)\) for the arithmetic mean of the oriented
fine-bond sums along the six coordinate-order paths from \(u\) to \(v\).
Set
\[
 \begin{aligned}
 \varphi_A(z,c)&=\exp\!\left(ie_k\eta
                    \sum_{j<k}\tau A(c_{j+1},c_j)\right),&
 (Q_k(A)\xi)(z)&=\eta^3\sum_{c\in\Delta_z}\varphi_A(z,c)\xi(c),\\
 (Q_k^T(-A)v)(c)&=\varphi_A(z,c)^{-1}v(z).&&
 \end{aligned}
\]
For \(z'\in L\Z^3/\ell\Z^3\), put
\(B(z')=\{z\in\Z^3:|z-z'|_\infty<L/2\}\), modulo \(\ell\Z^3\).
The next average and its algebraic reverse are
\[
 \begin{aligned}
 U_A(z',z)&=\exp(ie_k\eta\tau A(z',z)),\\
 (Q(A)v)(z')&=L^{-3}\sum_{z\in B(z')}U_A(z',z)v(z),&
 (Q^T(-A)v')(z)&=U_A(z',z)^{-1}v'(z').
 \end{aligned}
\]
The target pairing has measure \(L^3\). These maps act componentwise
in spin. For \(k\ge1\), define
\[
 \Pi_1(A)=JQ_k^T(-A)Q_k(A)J^*,\qquad
 \Pi_L(A)=JQ_k^T(-A)Q^T(-A)Q(A)Q_k(A)J^*.
\]
At \(k=0\) use only \(\Pi_L(A)=JQ^T(-A)Q(A)J^*\).
Let \(b=1\), \(b_L=L^{-1}\), \(b_-=1-L^{-1}\), and
\[
 b_k=\frac{1-L^{-1}}{1-L^{-k}}\quad(k\ge1),\qquad
 \alpha_y=\frac{b_ki\gamma_5y}{b_k+i\gamma_5y},\qquad
 \beta_y=\frac{b_k^2b_L}{(b_k+b_L+i\gamma_5y)(b_k+i\gamma_5y)}.
\]
At \(y=\infty\), put \(\alpha_\infty=b_k\), \(\beta_\infty=0\).
For \(0<\mu\le1\), the penalized operators are
\begin{equation}\label{eq:intro-penalized-operator}
 F_y(A)=
 \begin{cases}
 D_A+\mu+\alpha_y\Pi_1(A)+\beta_y\Pi_L(A),
      & k\ge1,\ y\in\R\cup\{\infty\},\\
 D_A+\mu+b_L\Pi_L(A)+i\gamma_5y,
      & k=0,\ y\in\R.
 \end{cases}
\end{equation}
The data \(k,\mu,\ell,e_k\) are independent.

Let \(M\ge L^2\) be a power of \(L\), with partition \(M\)-cubes
centered in \(M\Z^3/\ell\Z^3\), and assume \(\ell/(LM)\) is odd
and at least \(13\). A connection belongs to
\(\widetilde{\mathcal G}_k\) if \(A=A_0+A_1\), \(A_0\) is real,
\(\|A_1\|_\infty<e_k^{-1+\varepsilon}\), and on every concentric
\(11M\) enlargement of a partition cube,
\[
 A_0=B_0+d\lambda,\qquad B_0,\lambda\ \hbox{real},\qquad
 \|B_0\|_\infty<e_k^{-1+\varepsilon},\qquad
 d\lambda(c,c')=\eta^{-1}(\lambda(c')-\lambda(c)).
\]
The local gauges may be chosen independently.

Distances are periodic sup distances. For distinct chamber sites set
\(d_X(x,x')=\max(d(c,c'),\eta)\), and, on pairs with \(0<d_X<1\), put
\[
 U_A(x,x')=\exp(ie_k\eta\tau A(c,c')),\qquad
 (\delta_{\alpha,A}v)(x,x')
 =\frac{U_A(x,x')v(x')-v(x)}
        {d_X(x,x')^\alpha}.
\]
Thus distinct labels over one photon site are included when \(k\ge1\).
For an operator \(T\), define
\[
 N_\gamma(T)=\sup_z\sum_{z'}e^{\gamma d(z,z')}
       \|\mathbf1_{\Delta_z}T\mathbf1_{\Delta_{z'}}\|_{\infty\to\infty}.
\]
The sup norm uses Euclidean spinor norms; for pair outputs the left
indicator restricts the first point. At level zero,
\(N_\gamma^{\mathrm{ker}}\) replaces the localized operator norm by
the supremum of the normalized chamber kernel on the two clusters.

\begin{theorem}\label{thm:chamber-locality}
There are \(a'>0\) and \(0<\gamma\le(6L)^{-1}\), depending only on
\(\rho,L,\alpha\), with the following property. Set
\(C_V=2\exp(1)(42+6b)\) and
\(q(p)=54\pi C_Va'p/M\). Assume
\begin{equation}\label{eq:main-locality-conditions}
 M>4L,\qquad q(1)\le\tfrac12,\qquad
 0<e_k\le e_0,\qquad
 6Le_0^\varepsilon\le1,\qquad 4C_Va'e_0^\varepsilon\le1.
\end{equation}
For every \(A\in\widetilde{\mathcal G}_k\), all the operators
\eqref{eq:intro-penalized-operator} are invertible. Their inverses
\(S_y(A)\) are analytic in the connection and satisfy
\[
 N_\gamma(S_y(A))\le32a'w_k(y),\qquad
 N_\gamma(\delta_{\alpha,A}S_y(A))\le304a',\qquad
 w_k(y)=
 \begin{cases}1,&k\ge1,\\(1+|y|)^{-1},&k=0.\end{cases}
\]
At level zero the first bound also holds in
\(N_\gamma^{\mathrm{ker}}\); there are no short pairs at that level.
The constants are independent of the level, mesh, volume,
permitted connection, reference mass \(\mu\in(0,1]\), and resolvent
parameter in its stated range.
\end{theorem}

We develop the chamber block arguments of
\cite[Sections~9 and~11]{QED3Manuscript}. The internal chamber graph
has a positive gap on fields with zero cluster mean. Eliminating those
components gives a four-component Schur symbol with the Dirac
singularity at small momentum. Its Fourier bounds supply local
regularity uniformly as the mesh tends to zero. Averaging penalties
control the remaining block means. Local real gauges then reduce the
connection to a small perturbation, and two local inverse identities
give convergent right and left walk expansions.

The walk supports also localize the coarse covariance and determinant.
We first attach parameters to the entire support of each walk, so
that setting parameters to zero separates components. Exact resolvent
identities pass from the fine inverse to the coarse inverse.
Weighted Berezin integration then gives the block formula in
\cref{thm:fermion-block-integration}. A resolvent integral expresses
its determinant as a sum of local analytic terms. For the additional
small-field difference estimate, long polymers supply the gain by
decay; on short polymers a real gauge gives an enlarged complex disk,
and Cauchy's formula gives the gain. The necessary nonwrapping
geometry is proved in \cref{app:polymer-geometry}.

The separate free estimates below hold for arbitrary unitary period
characters. The integration and transpose identities use the
antiperiodic convention fixed here. Their spin matrices are
\[
 \mathsf C=i\gamma_2,\qquad
 \varkappa=\gamma_5\mathsf C=\sigma_2\otimes\sigma_2,
 \qquad
 \mathsf C^T=\mathsf C^{-1}=-\mathsf C,\quad
 \varkappa^T=\varkappa=\varkappa^{-1}.
\]
Both matrices are real. Their displayed forms give the transpose and
complex-conjugation identities used in the determinant calculation.
The connection is an input to these estimates. For a connection
obtained from Maxwell reconstruction, the local gauge and fluctuation
domain estimates are supplied separately by
\cite[Theorems~A and~B]{MaxwellBlockManuscript}.
\section{The free chamber inverse}\label{sec:propagators}

The chamber estimates of \cite[Section~9]{QED3Manuscript} give bounds which
are uniform in the mesh, the volume, and the running mass. There are two
sources of coercivity. The internal chamber graph controls the part of a
spinor whose mean over a cluster vanishes. The averaging penalty controls
the remaining four components. The first of these bounds also identifies
the singular part of the free propagator. This permits a local regularity
estimate with constants independent of the mesh, which is needed when
the inverses on small tori are combined into an inverse on the full torus.

\subsection{The cluster decomposition}\label{subsec:cluster-decomposition}

We first allow an arbitrary mesh $a>0$, using the chamber stencil of
\cref{sec:introduction} with $\eta$ replaced by $a$. Put
$\widehat d_i(c,g)=-gn_i$ and
\[
 M_i(c,g)=\frac{w_i}{2}(\rho-\gamma\cdot\widehat d_i(c,g)).
\]

The discussion applies on the covering lattice and on a torus of side
$R\in a\N$, with any unitary character
$f(c+Re_\mu,g)=-\exp(i\theta_\mu)f(c,g)$. Pairings on the torus are
formed over a period cell. The notation $d(c,c')$ denotes the sup
distance modulo $R\Z^3$, or the covering sup distance when $R=\infty$.
They use the measure
\[
 \int\dd x\,v(x)=\frac{a^3}{48}\sum_c\sum_gv(c,g),
 \qquad
 \|f\|_2^2=\int\dd x\,|f(x)|_2^2.
\]
On four-component spinors depending only on $c$, the measure is
$a^3\sum_c$. With the cluster maps $J,J^*$ of the introduction,
set $P=JJ^*$ and $Q=1-P$. The decomposition $f=J\xi+h$, with
$\xi=J^*f$ and $h=Qf$, is orthogonal. In particular,
\begin{equation}\label{eq:cluster-orthogonal}
 \|f\|_2^2=\|\xi\|_2^2+\|h\|_2^2.
\end{equation}
The normalization by $48$ makes $J$ an isometry. Kernels on the chamber
lattice are taken with respect to $\dd x$, so the kernel of the identity
is $48a^{-3}\mathbf1[x=x']1_4$.

Let $A$ be a connection and write
$u_A(c,c')=\exp(ie_a aA(c,c'))$ on a photon bond. The charge $e_a$ in
this notation is the charge at the level under consideration. The
phases on the internal bonds are one. Put
\[
 R_i(A)f(x)=U_A(x,x^{(i)})f(x^{(i)}),\qquad
 \mathcal L_A=\sum_{i=0}^3w_i(1-R_i(A)).
\]
The massless chamber operator is
\begin{equation}\label{eq:propagator-dirac}
 D_Af(x)=-a^{-1}\sum_{i=0}^3M_i(x)
       \big(U_A(x,x^{(i)})f(x^{(i)})-f(x)\big).
\end{equation}
At level $k$, $a=\eta_k$ and $e_a=e_k$. The reference mass is
added separately as $D_A+\mu$.
For real $A$, the involutions $R_i(A)$ are unitary and self-adjoint,
and the bond expansion gives
\begin{equation}\label{eq:dirac-realpart}
 \operatorname{Re}(f,D_Af)=\frac{\rho}{2a}(f,\mathcal L_Af),
 \qquad D_A^*=\gamma_5D_A\gamma_5,
 \qquad \gamma_5=\begin{pmatrix}0&1_2\\1_2&0\end{pmatrix}.
\end{equation}
Here $\gamma_5$ is Hermitian, $\gamma_5^2=1$, and it anticommutes with
each $\gamma_\mu$. Each unoriented bond contributes
$a^3w_i|U_A(x,x^{(i)})f(x^{(i)})-f(x)|_2^2/48$ to the energy.

On one cluster let
\[
 \mathcal L_{\mathrm{int}}=\sum_{i=0}^2w_i(1-R_i),
 \qquad R_iv(g)=v(gs_i).
\]
Its kernel consists of the constants: the three reflections
$s_0,s_1,s_2$ generate $O_{\mathrm{cub}}$. Denote its least positive
eigenvalue by $\lambda_{\mathrm{int}}$ and set
\[
 g_{\mathrm h}=\frac{\rho\lambda_{\mathrm{int}}}{2},\qquad
 c_* =\frac{\min(w_3,\lambda_{\mathrm{int}})}{18},\qquad
 c_{\mathrm{lh}}=\sqrt3(1+\rho).
\]
No numerical evaluation of $\lambda_{\mathrm{int}}$ is needed.

\begin{proposition}\label{prop:cluster-energy}
For real $A$ and $\xi=J^*f$,
\begin{equation}\label{eq:cluster-energy}
 (f,\mathcal L_Af)\ge\lambda_{\mathrm{int}}\|Qf\|_2^2,
 \qquad
 a^{-2}(f,\mathcal L_Af)\ge c_*\|\nabla_A\xi\|_2^2,
\end{equation}
where
$\nabla_{A,\mu}\xi(c)=a^{-1}(u_A(c,c+ae_\mu)\xi(c+ae_\mu)-\xi(c))$.
The same inequalities hold on a finite set $E$ of photon sites, with the
energy restricted to bonds whose endpoints lie over $E$ and the photon
differences restricted to bonds in $E$. Moreover,
\begin{equation}\label{eq:offdiagonal-bound}
 \|QD_AJ\xi\|_2\le c_{\mathrm{lh}}\|\nabla_A\xi\|_2,
 \qquad
 |(\xi,J^*D_AQf)|\le
 c_{\mathrm{lh}}\|\nabla_A\xi\|_2\|f\|_2.
\end{equation}
On $\operatorname{ran}Q$, for $t>-g_{\mathrm h}/a$,
\begin{equation}\label{eq:heavy-inverse}
 \|(QD_AQ+t)^{-1}\|_{2\to2}\le(t+g_{\mathrm h}/a)^{-1}.
\end{equation}
All constants are independent of $a,R$, the character, and the real
connection.
\end{proposition}

\begin{proof}
The internal bonds over $c$ give
$a^3(f(c,\cdot),\mathcal L_{\mathrm{int}}f(c,\cdot))$ and hence at
least $\lambda_{\mathrm{int}}a^3\|Qf(c,\cdot)\|^2$. Summing over
$c\in E$ proves the first bound in \eqref{eq:cluster-energy}.

For the second bound, consider a photon bond $c'=c+av$. The eight
labels with $ge_1=v$ correspond to the eight chamber bonds crossing it.
Write $\mathcal O_v(c)$ for this octet. Since their endpoints form
$\mathcal O_{-v}(c')$, the difference of the cluster means is the sum
of the mean of the eight bond differences and two means of $Qf$:
\begin{align*}
 u_A(c,c')\xi(c')-\xi(c)
 ={}&\frac18\sum_{x\in\mathcal O_v(c)}
       \big(U_A(x,x^{(3)})f(x^{(3)})-f(x)\big)\\
 &-\frac{u_A(c,c')}8\sum_{x'\in\mathcal O_{-v}(c')}Qf(x')
   +\frac18\sum_{x\in\mathcal O_v(c)}Qf(x).
\end{align*}
The square is bounded by $3/8$ times the sum of the squared terms in
these three octets. Distinct photon bonds use distinct chamber bonds;
each octet of cluster remainders occurs at most once. Consequently,
with $\mathcal E_3$ and $\mathcal E_{\mathrm{int}}$ denoting the two
parts of the energy,
\[
 a^2\|\nabla_A\xi\|_{2,E}^2
 \le18\big(w_3^{-1}\mathcal E_3+\|Qf\|_{2,E}^2\big)
 \le18\big(w_3^{-1}\mathcal E_3+
                \lambda_{\mathrm{int}}^{-1}\mathcal E_{\mathrm{int}}\big).
\]
This proves the local bound and therefore the global one.

The internal hops vanish on $J\xi$. Thus
\[
 (D_AJ\xi)(c,g)=-3(\rho-\gamma\cdot ge_1)
                   \nabla_{A,ge_1}\xi(c).
\]
Each of the six directions occurs eight times. The norm of the matrix
in this formula is $3(1+\rho)$, and the two directed differences on a
photon bond have equal norms. It follows that
$\|D_AJ\xi\|_2^2\le3(1+\rho)^2\|\nabla_A\xi\|_2^2$.
Projection gives the first bound in \eqref{eq:offdiagonal-bound}; the
second follows from \eqref{eq:dirac-realpart}, since $J$ and $Q$ commute
with $\gamma_5$.

For $Qf=f$, \eqref{eq:dirac-realpart} gives
$\operatorname{Re}(f,(QD_AQ+t)f)\ge(t+g_{\mathrm h}/a)\|f\|_2^2$.
The same lower bound holds for the adjoint. The operator has closed
range and its adjoint has zero kernel, so it is onto, also on the
infinite lattice. This proves \eqref{eq:heavy-inverse}.
\end{proof}

\subsection{The free symbol and its Schur complement}
\label{subsec:schur-symbol}

Use lattice units for the moment. Let $\mathcal D_m(p)$ be the symbol
of $D_0+m$, $0\le m\le1$, on
$V=\ell^2(O_{\mathrm{cub}};\C^4)$ with measure $48^{-1}\sum_g$:
\[
 (\mathcal D_m(p)v)(g)=m v(g)
 +\sum_{i=0}^2M_i(g)(v(g)-v(gs_i))
 +M_3(g)(v(g)-\exp(ip\cdot ge_1)v(gr_{e_1})).
\]
Let $D_{\mathrm{ll}},D_{\mathrm{lh}},D_{\mathrm{hl}},D_{\mathrm{hh}}$
be its blocks in $V=\operatorname{ran}J\oplus\operatorname{ran}Q$.
The mass is included in the diagonal blocks. Set
\[
 H_m(p)=D_{\mathrm{hh}}(p)^{-1},\qquad
 B_m(p)=D_{\mathrm{ll}}(p)-D_{\mathrm{lh}}(p)H_m(p)D_{\mathrm{hl}}(p),
 \qquad t(p)=\sum_\mu(1-\cos p_\mu).
\]
The compression of the full operator is the Wilson operator:
\[
 D_{\mathrm{ll}}(p)=m+i\gamma\cdot\sin p+\rho t(p).
\]
Its Schur complement also contains the contribution from the heavy
part. In particular, adding a mass to the chamber operator changes
$H_m$ as well as $D_{\mathrm{ll}}$.

\begin{proposition}\label{prop:schur-symbol}
There are positive constants $\theta_0,c_{\mathrm e},\kappa_{\mathrm{strip}}$, depending
on $\rho$ only, such that $H_m$ and $B_m$ are holomorphic in a
neighborhood of $\{|\operatorname{Im}p_\mu|\le\theta_0\}$ and
$2\pi\Z^3$-periodic. On the real cube $[-\pi,\pi]^3$,
\begin{align}
 \|H_m(p)\|&\le(g_{\mathrm h}+m)^{-1},\label{eq:symbol-heavy}\\
 \|D_{\mathrm{lh}}(p)\|+\|D_{\mathrm{hl}}(p)\|
   &\le2c_{\mathrm{lh}}(2t(p))^{1/2},\label{eq:symbol-cross}\\
 B_m(p)&=m+i\gamma\cdot\sin p+E_m(p),\qquad
 \|E_m(p)\|\le K_\rho t(p),\label{eq:schur-expansion}\\
 \operatorname{Re}(v,B_m(p)v)&\ge(m+\rho c_*t(p))|v|_2^2,
                 \label{eq:schur-realpart}\\
 |B_m(p)v|_2&\ge c_{\mathrm e}(|p|_2^2+m^2)^{1/2}|v|_2,
                 \label{eq:schur-ellipticity}
\end{align}
where $K_\rho=\rho+2c_{\mathrm{lh}}^2/g_{\mathrm h}$. For every
multi-index $\nu$,
\begin{equation}\label{eq:schur-inverse-derivatives}
 \|\partial_p^\nu B_m(p)^{-1}\|
       \le C_\nu(\rho)(|p|_2^2+m^2)^{-(1+|\nu|)/2}.
\end{equation}
If $m>0$, the inverse exists also on
$|\operatorname{Im}p_\mu|\le\kappa_{\mathrm{strip}}m$, where the same bound holds
with $\operatorname{Re}p$ on the right, for
$\operatorname{Re}p\in[-\pi,\pi]^3$. At $m=0$ the real derivative
bound is asserted away from $p=0$.
\end{proposition}

\begin{proof}
Only the fourth hop depends on $p$. For $p=p'+ip''$,
\[
 \|\mathcal D_m(p)-\mathcal D_m(p')\|
       \le3(1+\rho)(\exp(|p''|_\infty)-1).
\]
Choose $\theta_0$ so that the right side is $g_{\mathrm h}/2$ when
$|p''|_\infty=\theta_0$. The heavy gap then gives
$\operatorname{Re}D_{\mathrm{hh}}(p)\ge m+g_{\mathrm h}/2$ on this
strip. Invertibility holds on a slightly larger strip, which proves
holomorphy up to the closed strip. For real $p$, the off-diagonal bounds
in \cref{prop:cluster-energy} give
$\|D_{\mathrm{lh}}\|,\|D_{\mathrm{hl}}\|
\le c_{\mathrm{lh}}(2t(p))^{1/2}$. This proves
\eqref{eq:schur-expansion}.

To retain the real-part bound after eliminating the heavy part, put
$u=Jv-H_mD_{\mathrm{hl}}v$. Then
$J^*u=v$, $Q\mathcal D_m u=0$, and
$\mathcal D_m u=JB_mv$. Therefore
\[
 \operatorname{Re}(v,B_mv)
   =\operatorname{Re}(u,\mathcal D_m u)
   =m\|u\|^2+\frac\rho2(u,\mathcal L_0(p)u)
   \ge(m+\rho c_*t(p))|v|_2^2.
\]
This proves \eqref{eq:schur-realpart}.

The ellipticity estimate uses the Dirac part near zero and the real
part away from zero. Put $t_0=(4K_\rho^2)^{-1}$. One has
$K_\rho\ge48/\lambda_{\mathrm{int}}>1$, since
$\lambda_{\mathrm{int}}\le12+24\sqrt2<48$.
Hence $t_0\le1/4$. If $t(p)\le t_0$, then
$|\sin p|_2^2\ge t(p)$ and
$\|E_m(p)\|\le\tfrac12t(p)^{1/2}$. Since
\[
 |(m+i\gamma\cdot\sin p)v|_2^2=(m^2+|\sin p|_2^2)|v|_2^2,
\]
it follows that
$|B_m(p)v|_2\ge\tfrac12(m^2+t(p))^{1/2}|v|_2$.
On $[-\pi,\pi]^3$, $t(p)\ge2|p|_2^2/\pi^2$. This proves
\eqref{eq:schur-ellipticity} in this region. If $t(p)\ge t_0$, then
\eqref{eq:schur-realpart} gives
$|B_m(p)v|_2\ge(m+\rho c_*t_0)|v|_2$.
The inequality $(|p|_2^2+m^2)^{1/2}\le m+\pi\sqrt3$ closes the
estimate with, for example,
\[
 c_{\mathrm e}=\min\left(0.22,\frac{\rho c_*t_0}{\pi\sqrt3}\right).
\]

The strip bounds give uniform bounds on the derivatives of $B_m$ by
Cauchy's formula. Differentiating $B_mB_m^{-1}=1$ and inducting on
$|\nu|$ proves \eqref{eq:schur-inverse-derivatives}. For the complex
strip, let $K_s$ bound $B_m$ on the fixed strip. On half that strip,
$\|\partial_\mu B_m\|\le2K_s/\theta_0$. Choose
\[
 \kappa_{\mathrm{strip}}=\min\left(\frac{\theta_0}{2},
                \frac{c_{\mathrm e}\theta_0}{12K_s}\right).
\]
Integration along the segment from $p'$ to $p'+ip''$ shows that the
change in $B_m$ is at most $c_{\mathrm e}m/2$. The lower bound is thus
at least $c_{\mathrm e}(|p'|_2^2+m^2)^{1/2}/2$. Differentiating the
inverse identity once more gives the asserted complex derivative bounds.
\end{proof}

\subsection{Free kernels and coercivity}\label{subsec:free-kernels}

Solving the heavy equation gives the inverse symbol
\begin{equation}\label{eq:inverse-symbol-blocks}
 \mathcal D_m^{-1}=
 \begin{pmatrix}
 B_m^{-1}&-B_m^{-1}D_{\mathrm{lh}}H_m\\
 -H_mD_{\mathrm{hl}}B_m^{-1}&
 H_m+H_mD_{\mathrm{hl}}B_m^{-1}D_{\mathrm{lh}}H_m
 \end{pmatrix}.
\end{equation}
The first term is independent of the cluster labels after insertion by
$J$ and $J^*$. Every other term has a weaker singularity at $p=0$:
$D_{\mathrm{lh}}$ and $D_{\mathrm{hl}}$ vanish there, whereas $H_m$
and all its derivatives are uniformly bounded.

\begin{proposition}\label{prop:free-chamber-kernel}
On the unit covering lattice let
\[
 G_m(c,g;c',g')=\int_{[-\pi,\pi]^3}\frac{\dd p}{(2\pi)^3}
       \exp(ip\cdot(c-c'))[\mathcal D_m(p)^{-1}](g,g').
\]
The bracket is the kernel on $V$ relative to the label measure
$48^{-1}\sum_g$, consistently with the chamber kernel convention.
It is well-defined also at $m=0$. It decomposes as
$G_m(x,x')=G_m^{\mathrm l}(c-c')+G_m^{\mathrm h}(x,x')$, where
\begin{align}
 |G_m^{\mathrm l}(n)|_2&\le C_G(1+|n|_\infty)^{-2},\label{eq:free-light}\\
 |G_m^{\mathrm l}(n+v)-G_m^{\mathrm l}(n)|_2
       &\le C_G(1+|n|_\infty)^{-3},\quad v\in\{\pm e_\mu\},\\
 |G_m^{\mathrm h}(x,x')|_2&\le C_G(1+|c-c'|_\infty)^{-3}.
\end{align}
Consequently $|G_m(x,x')|_2\le2C_G(1+|c-c'|_\infty)^{-2}$ and
\begin{equation}\label{eq:free-bond-difference}
 |G_m(x^{(i)},x')-G_m(x,x')|_2
       \le20C_G(1+|c-c'|_\infty)^{-3}.
\end{equation}
For $m>0$ these bounds have the additional factor
$\exp(-\kappa_{\mathrm{strip}}m|c-c'|_\infty)$. The constants are uniform for
$0\le m\le1$.
\end{proposition}

\begin{proof}
The Fourier estimate used here can be proved directly. If a periodic
symbol $T$ satisfies
$|\partial^\nu T(p)|\le A|p|_2^{-q-|\nu|}$, where $q<3$ and
$|\nu|\le N>3-q$, then its Fourier coefficients are bounded by
$C A(1+|n|_\infty)^{-(3-q)}$. For $n\ne0$, put
$r=|n|_\infty^{-1}$ and split the integral using a cutoff equal to
one for $|p|_2\le r$ and zero for $|p|_2\ge2r$. The inner integral
is bounded by $CA r^{3-q}$. On the complement integrate $N$ times in
a coordinate with $|n_\mu|=|n|_\infty$. The differentiated integrand
has integral at most $CA r^{3-q-N}$; the prefactor is $r^N$.
Derivatives of the cutoff are supported in $r\le|p|_2\le2r$ and have
the same bound. Boundary terms cancel by periodicity.

Apply this estimate with $(q,N)=(1,3)$ to $B_m^{-1}$. Multiplication
by $\exp(ip\cdot v)-1$ lowers its singularity to order zero, for which
$(q,N)=(0,4)$ applies. In \eqref{eq:inverse-symbol-blocks}, the two
off-diagonal terms have order zero and the term with two off-diagonal
factors has order $-1$. The latter can also be bounded by the
order-zero estimate. The remaining term $H_m$ is holomorphic on the
fixed strip, and its Fourier coefficients decay exponentially.
The kernel normalization costs a factor $48$ for operator-valued
symbols on $V$; this is absorbed in $C_G$.

For $m>0$, shift one coordinate of the contour by
$i\kappa_{\mathrm{strip}}m\operatorname{sgn}(n_\mu)$, where
$|n_\mu|=|n|_\infty$. The complex derivative bounds of
\cref{prop:schur-symbol} give the same Fourier estimates
on the shifted contour. This supplies the exponential factor.
The fixed-strip bound for $H_m$ is stronger than the required one.

An internal hop leaves $c$ fixed, so its difference cancels
$G_m^{\mathrm l}$. An external hop replaces $c$ by $c+ge_1$; use the
light difference bound and the two heavy bounds at these points.
The inequalities $1+|n+v|_\infty\ge(1+|n|_\infty)/2$ and
$\exp(\kappa_{\mathrm{strip}}m)<\sqrt3$ give the factor $20$ in
\eqref{eq:free-bond-difference}.
\end{proof}

The Green kernel at mesh $a$, for $0\le\mu a\le1$, is
$a^{-2}G_{\mu a}(x/a,x'/a)$. Thus
\begin{equation}\label{eq:mesh-free-kernel}
 |G(x,x')|_2\le2C_G(a+|c-c'|_\infty)^{-2},\qquad
 |G(x^{(i)},x')-G(x,x')|_2
       \le20C_Ga(a+|c-c'|_\infty)^{-3}.
\end{equation}
Fourier multiplication gives both inverse identities on finitely
supported functions. The factor $48$ in the identity kernel is
consistent with the measure $\dd x$ above.

For distinct sites put
$d_X(x,\widetilde x)=\max(|c-\widetilde c|_\infty,a)$.
The chamber graph has paths from $x$ to $\widetilde x$ of at most
$39d_X(x,\widetilde x)/a$ bonds, with photon sites in the coordinate
box spanned by $c,\widetilde c$. Indeed, the generators
$s_0,s_1,s_2$ are the simple reflections of the signed permutation
group, with nine positive roots. Their Coxeter length is at most nine
\cite[Theorem~22.14]{EtingofLieNotes}, so any two labels in one
cluster are joined by at most nine internal bonds. A monotone photon
path has at most $3|c-\widetilde c|_\infty/a$ steps. Before each step
move to a label with the required external direction, and after the
last step move to the final label. This uses at most
$10\cdot3|c-\widetilde c|_\infty/a+9\le39d_X/a$ bonds.
If $d_X$ is less than half the distance to $c'$, sum
\eqref{eq:mesh-free-kernel} along such a path. In the other case use
the two size bounds. The resulting interpolation gives, for
$0\le\alpha\le1$,
\begin{equation}\label{eq:free-holder-kernel}
 |G(x,x')-G(\widetilde x,x')|_2
 \le6240C_G d_X(x,\widetilde x)^\alpha
 \left((a+|c-c'|_\infty)^{-2-\alpha}
       +(a+|\widetilde c-c'|_\infty)^{-2-\alpha}\right).
\end{equation}

The following coercivity estimate is separate from the real-part bound.
It is needed to recover mesh-independent control of the block remainder.

\begin{proposition}\label{prop:free-coercivity}
There is $c_b=c_b(\rho)>0$ such that
\begin{equation}\label{eq:free-coercivity}
 \|D_0f\|_2^2\ge c_ba^{-2}(f,\mathcal L_0f).
\end{equation}
For a cube $E$ of side $\ell_E$, let $P_E$ be the mean over
$E\times O_{\mathrm{cub}}$. Then
\begin{equation}\label{eq:chamber-poincare}
 \|f-P_Ef\|_{2,E}^2\le C_P\ell_E^2a^{-2}(f,\mathcal L_0f)_E,
 \qquad C_P=\frac1{8c_*}+\frac1{\lambda_{\mathrm{int}}}.
\end{equation}
\end{proposition}

\begin{proof}
The Schur estimate at zero mass implies
$|B_0(p)v|\ge c_{\mathrm l}t(p)^{1/2}|v|$, with
$c_{\mathrm l}=\rho c_*/(2K_\rho)$: use the Dirac part for
$t(p)\le t_0$ and the real part for $t(p)\ge t_0$.
Write $u=Jv+h$ in a fibre and $f=\mathcal D_0(p)u$. The heavy equation
and then the light equation give
\[
 h=H_0(f_{\mathrm h}-D_{\mathrm{hl}}v),\qquad
 B_0v=f_{\mathrm l}-D_{\mathrm{lh}}H_0f_{\mathrm h}.
\]
Since $t(p)\le6$, these imply
$(2t(p))^{1/2}|v|\le A_1\|f\|$ and $\|h\|\le A_2\|f\|$, where
\[
 A_1=\frac{\sqrt2}{c_{\mathrm l}}
         \left(1+\frac{6(1+\rho)}{g_{\mathrm h}}\right),\qquad
 A_2=g_{\mathrm h}^{-1}(1+c_{\mathrm{lh}}A_1).
\]
The bond decomposition also gives
\[
 (u,\mathcal L_0(p)u)\le6t(p)|v|^2+(30+24\sqrt2)\|h\|^2.
\]
For the internal bonds the light part cancels, leaving at most
$2(w_0+w_1+w_2)\|h\|^2$. In an external bond difference, split off
the photon difference of $v$ and the two heavy endpoint values, and
bound the square by three times their squared norms. The octet count
then gives $6t(p)|v|^2+3w_3\|h\|^2$. This proves the displayed
upper bound.
This yields \eqref{eq:free-coercivity} in each fibre with
$c_b=(3A_1^2+(30+24\sqrt2)A_2^2)^{-1}$.
Fourier decomposition gives the assertion for all characters on the
torus and on the covering lattice, restoring the factor $a^{-2}$.

For \eqref{eq:chamber-poincare}, first apply the scalar lattice
Poincar\'e bound
$\|\xi-P_E\xi\|_{2,E}^2\le\ell_E^2\|\nabla_0\xi\|_{2,E}^2/8$.
In one dimension this follows by writing each difference as a sum of
nearest-neighbor differences; the coefficient of a bond at position
$j$ is $j(n-j)/2\le n^2/8$. In three dimensions use the three commuting
line averages, whose successive remainders are orthogonal.
Now $f-P_Ef=J(\xi-P_E\xi)+Qf$. Its two parts are orthogonal, and the
local bounds \eqref{eq:cluster-energy} give
\eqref{eq:chamber-poincare}. The heavy term uses $\ell_E\ge a$.
\end{proof}

\section{Local inverses and walk expansions}\label{sec:local-walks}

\subsection{Averaging penalties and the inverse at zero field}
\label{subsec:blocked-zero-inverse}

Fix an odd $L\ge5$ and put $a=\eta_k=L^{-k}$. Write $e_k$ for the
charge $e_a$ at this mesh. The level data $k,\mu,R,e_k$ are independent,
with $0<e_k\le e_0$.
Unit blocks have side
one, and the next blocks have side $L$. In this and the following
subsections a finite period $R$ is a positive integer multiple of $L$.
For $z\in\Z^3$ the unit block is
$\Delta_z=\{c\in a\Z^3:|c-z|_\infty<1/2\}$.
For $z'\in L\Z^3$ the $L$-block is
$E_{z'}=\{c\in a\Z^3:|c-z'|_\infty<L/2\}$.
These partitions are taken modulo $R\Z^3$ on the torus. Since $L$
is odd, no fine site lies on a block boundary. Their means, including
the cluster mean, are denoted $\Pi_1$ and $\Pi_L$. At nonzero connection
the means use the gauge phases of \cref{sec:introduction}.
With the constants and coefficients defined there, the blocked operators
at zero connection are
\begin{equation}\label{eq:blocked-zero-operator}
 F_y(0)=D_0+\mu+\alpha_y\Pi_1+\beta_y\Pi_L\quad(k\ge1),
 \qquad
 F_y(0)=D_0+\mu+b_L\Pi_L+i\gamma_5y\quad(k=0).
\end{equation}
Here $0\le\mu\le1$. In the first formula $y\in\R\cup\{\infty\}$;
in the second $y\in\R$. We use $-\infty=\infty$ when an adjoint
identity reverses this parameter. The notation $F_y$ denotes an operator in this
section, not the curvature.

The penalty bounds follow by diagonalizing $\gamma_5$. On its eigenspace
of sign $\epsilon$, the two eigenvalues and their sum are
\[
 \frac{b_ki\epsilon y}{b_k+i\epsilon y},\quad
 \frac{b_k^2b_L}{(b_k+b_L+i\epsilon y)(b_k+i\epsilon y)},\quad
 \frac{b_k(b_L+i\epsilon y)}{b_k+b_L+i\epsilon y}.
\]
For the first eigenvalue $v_1$ and the sum $v_2$, direct calculation gives
\[
 |v_1|^2=b_k\operatorname{Re}v_1,\qquad
 |v_2|^2\le b_k\operatorname{Re}v_2,\qquad
 \operatorname{Re}v_2\ge\frac{b_kb_L}{b_k+b_L}\ge\frac{b_-}{L}.
\]
Also $|\alpha_y|\le b_k\le b$ and $|\beta_y|\le b_L$.
On an $L$-block the decomposition
$\operatorname{Pen}=\alpha_y(\Pi_1-\Pi_L)+(\alpha_y+\beta_y)\Pi_L$
has orthogonal ranges. The scalar inequalities on the two spin
eigenspaces therefore give
\begin{equation}\label{eq:penalty-estimates}
 \operatorname{Re}(f,\operatorname{Pen}f)_E
      \ge\frac{b_-}{L}\|P_Ef\|_{2,E}^2,
 \qquad
 \|\operatorname{Pen}f\|_{2,E}^2
      \le b\operatorname{Re}(f,\operatorname{Pen}f)_E.
\end{equation}
The kernel on $E$ is
$\alpha_y\mathbf1[c,c'\text{ in the same unit block}]+\beta_yL^{-3}$,
and its norm is at most $2b$.

\begin{proposition}\label{prop:zero-blocked-inverse}
There are $C_I,C_5,\gamma_5'>0$, depending on $L,\rho$ and a fixed
$\alpha\in(0,1)$ only, such that $F_y(0)$ is invertible and
\[
 \|F_y(0)^{-1}\|_{2\to2}\le LC_I\qquad(k\ge1).
\]
If $f$ is supported over one unit block $\Delta_{z'}$ and
$v=F_y(0)^{-1}f$, then for $x$ over $\Delta_z$ and
$0<d_X(x,\widetilde x)\le1$,
\begin{equation}\label{eq:zero-blocked-holder}
 |v(x)|_2+d_X(x,\widetilde x)^{-\alpha}|v(x)-v(\widetilde x)|_2
 \le C_5\exp(-\gamma_5'd(z,z'))\|f\|_\infty.
\end{equation}
This includes pairs in one cluster. The constants are uniform in
$k,R$, the character, $\mu$, and $y$.
\end{proposition}

\begin{proof}
First suppose $k\ge1$. Put $f=F_y(0)u$ and
$X=\operatorname{Re}(u,f)$. The real-part identity
and \eqref{eq:penalty-estimates} imply
\[
 \sum_E\|P_Eu\|_{2,E}^2\le\frac L{b_-}X,\qquad
 \|\operatorname{Pen}u\|_2^2\le bX,\qquad
 \|\mu u\|_2^2\le X.
\]
Therefore $\|D_0u\|_2\le\|f\|_2+(1+\sqrt b)X^{1/2}$.
Use \eqref{eq:free-coercivity} for $D_0u$ and
\eqref{eq:chamber-poincare} on each $L$-block. With
$K_1=\sqrt L+\sqrt{Lb}$ this gives
\[
 \|u\|_2^2\le L\left(\frac1{b_-}+\frac{2C_PK_1^2}{c_b}\right)
           \|u\|_2\|f\|_2
       +\frac{2C_P}{c_b}L^2\|f\|_2^2.
\]
Hence $\|u\|_2\le LC_I\|f\|_2$, where
$C_I=b_-^{-1}+2C_PK_1^2/c_b+(2C_P/c_b)^{1/2}$.
The identity $F_y(0)^*=\gamma_5F_{-y}(0)\gamma_5$ gives the same
lower bound for the adjoint. Thus $F_y(0)$ has closed, dense range
and is onto, including on the infinite lattice.

To obtain decay, let $\psi$ be a bounded real periodic function with
$|\psi(c)-\psi(c')|\le a_0|c-c'|_1/L$. The mass commutes with this
weight. The external hopping term gives
$\|\exp(\psi)[D_0,\exp(-\psi)]\|_{2\to2}
\le3(1+\rho)\exp(1)a_0/L$.
The penalty kernel is supported in one $L$-block, has norm at most
$2b$, and $|\psi(c)-\psi(c')|\le3a_0$ there. Its contribution is at
most $6L^4b\exp(1)a_0/L$. Put
\[
 c_w=\exp(1)(3(1+\rho)+6L^4b),\qquad
 a_0=\min(1/3,(2C_Ic_w)^{-1}).
\]
Applying the $L^2$ inverse bound and absorbing the commutator gives
\begin{equation}\label{eq:weighted-block-inverse}
 \|\exp(\psi)F_y(0)^{-1}f\|_2
       \le2LC_I\|\exp(\psi)f\|_2.
\end{equation}
Take truncated multiples of the $\ell^1$ distance, divided by $L$,
to the source $L$-block. Letting the truncation increase proves an
exponentially weighted $L^2$ bound.

The passage from $L^2$ to the sup norm uses the free kernel. On the
covering lattice choose $c^\circ\in a\Z^3$ and $r\ge8a$, and put
\[
 B=\{|c-c^\circ|_\infty<r\},\quad
 B^+=\{|c-c^\circ|_\infty<r+a\},\quad
 B'=\{|c-c^\circ|_\infty\le r/4\}.
\]
Suppose $u$ is defined over $B^+$ and $(D_0+\mu)u=h$ over $B$,
with $0\le\mu a\le1$. Both evaluation points below lie over $B'$,
and $C$ depends on the fixed $\alpha\in(0,1)$ and $\rho$.
Let $\chi$ equal one
on the concentric half cube and zero outside the cube, with Lipschitz
constant $2/r$. Then
\[
 \chi u=G(\chi h)+G([D_0,\chi]u).
\]
The commutator contains only the external hop. Its rows are supported
where $|c-c^\circ|_\infty>r/2-a$. It satisfies
$\|[D_0,\chi]u\|_1\le45r^{1/2}\|u\|_{2,B^+}$.
For a point of the quarter cube, this support is at distance at least
$r/4-a\ge r/8$. The bounds \eqref{eq:mesh-free-kernel} and
\eqref{eq:free-holder-kernel} consequently give
\begin{equation}\label{eq:interior-dirac-estimate}
 \begin{split}
 |u(x)|_2&\le C\big(r\|h\|_{\infty,B}+r^{-3/2}\|u\|_{2,B^+}\big),\\
 |u(x)-u(\widetilde x)|_2&\le
 C\left(\frac{d_X(x,\widetilde x)}r\right)^\alpha
       \big(r\|h\|_{\infty,B}+r^{-3/2}\|u\|_{2,B^+}\big).
 \end{split}
\end{equation}
For the first term, sum $(a+|c-c'|_\infty)^{-2}$ over the cube;
the result is $O(r)$. For the difference, the analogous sum with
exponent $2+\alpha$ is $O(r^{1-\alpha}/(1-\alpha))$.
The cluster average in $\dd x'$ contributes one to both sums.
For the commutator term, its support separation gives the size bound;
summing bond differences along the chamber path gives the difference
bound. These estimates explain why the constants do not grow when
$a\downarrow0$.

For $u=F_y(0)^{-1}f$, use
$h=f-\operatorname{Pen}u$ in \eqref{eq:interior-dirac-estimate}, with
$r=L/4$. The enlarged cube meets at most eight $L$-blocks. On each,
the penalty kernel and Cauchy--Schwarz give
$\|\operatorname{Pen}u\|_\infty\le2bL^{3/2}\|u\|_{2,E}$.
The weighted bound \eqref{eq:weighted-block-inverse} therefore bounds
both terms in \eqref{eq:interior-dirac-estimate} by
$CL\exp(-a_0d_1(x,E'))\|f\|_\infty$, where
$d_1(x,E')$ is the $\ell^1$ distance to the source $L$-block divided
by $L$. This applies for $k\ge2$.
For $k=1$, the pointwise $L^2$ bound costs only
$\sqrt{48}L^{3/2}$; $L$ is fixed. Larger Hölder separations are
controlled by the two sup bounds. Finally,
$d_1(x,E')\ge d(z,z')/L-3$, and the source unit block lies in $E'$.
This proves \eqref{eq:zero-blocked-holder} with
$\gamma_5'=a_0/L$. The argument uses only periodic moduli and so
applies to every unitary character.
\end{proof}

The first level has a different resolvent parameter. On an $L$-block
the real-part bound and Poincar\'e inequality imply
\[
 \operatorname{Re}(u,F_y(0)u)\ge\frac{c_r}{L}\|u\|_2^2,
 \qquad c_r=\min(b_-,\rho/(2C_PL)).
\]
The term $i\gamma_5y$ has zero real part. Conjugation by a weight with
Lipschitz constant $a_z$ changes the hopping and penalty terms by at
most $\exp(1)(6+3b)a_z$. Choose
\[
 a_z\le\min\bigl((3L)^{-1},c_r/(2\exp(1)L(6+3b))\bigr).
\]
The weighted inverse then has $L^2$ norm at most $2L/c_r$.
On the other hand its part other than $i\gamma_5y$ has norm at most
$60+b$. For $|y|\ge2(60+b)$, the inverse norm is at most $2/|y|$.
Combining the two ranges and converting to a kernel, with its factor
$48$, gives
\begin{equation}\label{eq:level-zero-kernel}
 |F_y(0)^{-1}(x,x')|_2\le C_{00}(1+|y|)^{-1}
                         \exp(-a_zd(c,c')).
\end{equation}
This factor in $y$ is needed when the determinant is represented by an
integral of resolvents.

\subsection{Small connections and local inverses}\label{subsec:local-inverses}

Use the covariant block means, short quotients, and row norms of
\cref{sec:introduction}. Replacing $D_0$ by $D_A$ and $\Pi_1,\Pi_L$
by their covariant versions in \eqref{eq:blocked-zero-operator}
gives $F_y(A)$. For real
$A$ these idempotents are orthogonal projections.
Inserting the block partition between two factors and
using the triangle inequality for $d$ proves
$N_\gamma(TT')\le N_\gamma(T)N_\gamma(T')$.
Enlarge $C_5,C_{00}$ to at least one, and choose
\[
 \gamma=\tfrac12\min(\gamma_5',a_z),\qquad
 a'=\max(C_5,C_{00})\sum_{n\in\Z^3}\exp(-\gamma|n|_\infty).
\]
The preceding estimates imply
\begin{equation}\label{eq:zero-row-bounds}
 N_\gamma(F_y(0)^{-1})\le a'w_k(y),\qquad
 N_\gamma(\delta_{\alpha,0}F_y(0)^{-1})\le a'.
\end{equation}
At $k=0$ the first bound holds also in $N_\gamma^{\mathrm{ker}}$.

\begin{proposition}\label{prop:small-connection-inverse}
Let $B$ be a complex connection with $e_k\|B\|_\infty\le\delta$, where
$3L\delta\le1$ and $2C_Va'\delta\le1$, with
$C_V=2\exp(1)(42+6b)$. Then $F_y(B)$ is invertible and
\begin{equation}\label{eq:small-row-bounds}
 N_\gamma(F_y(B)^{-1})\le2a'w_k(y),\qquad
 N_\gamma(\delta_{\alpha,B}F_y(B)^{-1})\le18a'.
\end{equation}
At $k=0$ the first bound holds also in $N_\gamma^{\mathrm{ker}}$.
\end{proposition}

\begin{proof}
Write $V=F_y(B)-F_y(0)$. The internal hops, the mass, and the
resolvent term at level zero cancel. For the external hop,
$|u_B-1|\le\eta_k\delta\exp(\eta_k\delta)$; its weighted row sum
is at most $42\exp(\gamma)\exp(\delta)\delta$. The two averaging
penalties contribute respectively
$3b\delta\exp(3\delta)$ and
$3b\delta\exp(\gamma(L-1))\exp(3L\delta)$. Indeed, the nested paths
have total length at most $3/2$, and the path to an $L$-block center
has length at most $3L/2$. Summing their normalized block kernels
removes the numbers of sites. The bounds $\gamma L\le\log2$ and
$3L\delta\le1$ give $N_\gamma(V)\le C_V\delta$.

Thus $N_\gamma(VF_y(0)^{-1})\le1/2$. The Neumann series gives
\[
 F_y(B)^{-1}=F_y(0)^{-1}
                  \sum_{n\ge0}(-VF_y(0)^{-1})^n
             =\sum_{n\ge0}(-F_y(0)^{-1}V)^nF_y(0)^{-1}.
\]
It proves the first bound in \eqref{eq:small-row-bounds}. The first
form gives $N_\gamma(\delta_{\alpha,0}F_y(B)^{-1})\le2a'$.
For distinct photon sites in a short pair,
$d_X^{-\alpha}|U_B-1|\le3\delta\exp(3\delta)\le1/4$;
for intracluster pairs the expression is zero. Evaluation at the
second point of a short pair uses at most $27$ neighboring blocks,
and costs at most $27\exp(\gamma)$. Its contribution to the covariant
quotient is at most $16a'$, proving the second bound. At level zero,
the second form of the Neumann series keeps the kernel norm in the
last factor, so the kernel estimate follows as well.
\end{proof}

Return to the side $\ell=R$ and the domain
$\widetilde{\mathcal G}_k$ of \cref{sec:introduction}.
On the concentric $3M$ enlargement $W_\square$ of an $M$-cube,
\begin{equation}\label{eq:local-small-gauge}
 A=B+d\lambda,\qquad e_k\|B\|_\infty<2e_k^\varepsilon,
 \qquad e_k\|\operatorname{Im}A\|_\infty<e_k^\varepsilon.
\end{equation}
This follows by restricting $B_0+A_1$ and $\lambda$ to that cube.
The domain is open: a sufficiently small complex perturbation can be
added to $A_1$. It is invariant under real gauge transformations.

For a lift $z_\square$ of the cube center put
$W_\square=\{|c-z_\square|_\infty<3M/2\}$ and
$W_\square^-=\{|c-z_\square|_\infty<3M/2-\eta_k\}$.
Embed $W_\square$ in the positive chamber of the periodic torus of
side $6M$, whose lattice is translated by $\eta_k(1,1,1)/2$.
Its level-$j$ block centers are
$L^j\eta_k(1,1,1)/2+L^j\eta_k\Z^3$ for $0\le j\le k+1$.
Translation by $L(1,1,1)/2$ identifies these blocks with the
unshifted nested blocks, because $L^{k+1-j}$ is odd. The zero-field
and small-connection estimates therefore apply on this torus with
the same constants.
Fold its photon coordinates across the planes at multiples of $3M$.
Adjacent sites fold either to adjacent sites in $W_\square$ or to
one site; in the latter case assign phase one. Pull back $A$ by this
fold to obtain $A^W$. The fold preserves the supremum norm and
commutes with $d$ on scalar functions. It also preserves the nested
blocks, since $3M/(L^j\eta_k)$ is odd for $j\le k+1$.

Let $\mathrm{ext}$ extend a function over $W_\square$ by zero in the
seven other chambers, and let $\mathrm{res}$ restrict to its embedded
copy. Labels are unchanged in these two maps. Define
\[
 S_\square(A)=\mathrm{res}\,F_y(A^W)^{-1}\mathrm{ext}.
\]
For $A=B+d\lambda$ as in \eqref{eq:local-small-gauge},
$A^W=B^W+d\lambda^W$, with $\lambda^W$ real. Gauge conjugation by
$\exp(ie_k\lambda^W)$ reduces the inverse to the small connection
$B^W$. If $6Le_0^\varepsilon\le1$ and
$4C_Va'e_0^\varepsilon\le1$, \cref{prop:small-connection-inverse}
applies to every $0<e_k\le e_0$. Restriction, extension, and unitary
gauge factors do not increase the localized norms. Thus
\begin{equation}\label{eq:local-inverse-bounds}
 N_\gamma(S_\square(A))\le2a'w_k(y),\qquad
 N_\gamma(\delta_{\alpha,A}S_\square(A))\le18a'.
\end{equation}
The second operator here is formed on the local cube before extension
to the original torus. It does not differentiate a discontinuous
extension by zero.

The local inverse has both identities needed for patching. If $f$ is
defined over $W_\square$, then
\begin{equation}\label{eq:local-right-identity}
 F_y(A)S_\square(A)f=f\quad\text{over }W_\square^-.
\end{equation}
If $f$ is supported over $W_\square^-$, then $F_y(A)f$ is supported
over $W_\square$ and
\begin{equation}\label{eq:local-left-identity}
 S_\square(A)F_y(A)f=f.
\end{equation}
For \eqref{eq:local-right-identity}, the rows at an interior site use
only its internal neighbors, its external neighbor, and its
$L$-block. These lie in $W_\square$, where the original and folded
operators agree. For \eqref{eq:local-left-identity}, the support
margin ensures that applying $F_y$ cannot leave $W_\square$.
After embedding, applying the folded operator to the extension of
$f$ agrees with extending $F_y(A)f$. Apply the inverse and restrict.
For an antiperiodic spinor on the original torus, use one lift and
then its signed translates; the enlarged cubes do not overlap their
translates. These arguments also show that $S_\square$ reads $A$
only in $W_\square$ and is analytic in those connection variables.

\subsection{The parametrix and the walk expansion}\label{subsec:propagator-walks}

Choose the product cutoff $h_\square$ with center $z_\square$ and
one-dimensional factor
\[
 h_0(t)=\begin{cases}
 1,&|t|\le1/4,\\
 \cos(\pi(2|t|-1/2)/2),&1/4\le|t|\le3/4,\\
 0,&|t|\ge3/4.
 \end{cases}
\]
On a lift of the cube set
\[
 h_\square(c)=\prod_{\mu=1}^3
 h_0\bigl((c_\mu-z_{\square,\mu})/M\bigr),\qquad
 h_\square(c,g)=h_\square(c).
\]
Extend this with period $R$ and set it to zero away from the periodic
images of the lifted cube neighborhood.
The cutoffs satisfy $\sum_\square h_\square^2=1$,
$\supp h_\square\subset\{|c-z_\square|_\infty<3M/4\}$, and
$|h_\square(c)-h_\square(c')|\le\pi|c-c'|_1/M$.
Put $C_\square=-[F_y(A),h_\square]$. Internal hops commute with
$h_\square$, and the penalty has range at most $L$. The phase bounds
and the same normalized sums used in the perturbation estimate give
\begin{equation}\label{eq:cutoff-commutator-bound}
 N_\gamma(C_\square)\le\pi C_V/M.
\end{equation}
The kernel is
$(h_\square(c)-h_\square(c'))F_y(A)(x,x')$. Its two photon sites lie
within $3M/4+L$ of $z_\square$. If $M>4L$, both lie in
$W_\square^-$; multiplication by $h_{\square'}$ can give a nonzero
operator only when $d(z_\square,z_{\square'})\le M$.

Using \eqref{eq:local-right-identity} and
\eqref{eq:local-left-identity} gives
\begin{align}
 F_yh_\square S_\square h_\square
       &=h_\square^2-C_\square S_\square h_\square,
                                   \label{eq:parametrix-right}\\
 h_\square S_\square h_\square F_y
       &=h_\square^2+h_\square S_\square C_\square.
                                   \label{eq:parametrix-left}
\end{align}
Let $S^{\mathrm p}=\sum_\square h_\square S_\square h_\square$ and
$R^{\mathrm p}=\sum_\square C_\square S_\square h_\square$.
At most $27$ cubes contribute to one row of $R^{\mathrm p}$.
Consequently,
\[
 N_\gamma(R^{\mathrm p})\le q_0:=\frac{54\pi C_Va'}M.
\]
Choose $M\ge108\pi C_Va'$. Equation~\eqref{eq:parametrix-right}
now gives a right inverse by a convergent series:
\begin{equation}\label{eq:global-parametrix-inverse}
 S_y(A)=F_y(A)^{-1}=S^{\mathrm p}\sum_{n\ge0}(R^{\mathrm p})^n.
\end{equation}
The finite-dimensional operator $F_y(A)$ is therefore invertible.
The left identity gives the second form with
$R^{\mathrm p\prime}=-\sum_\square h_\square S_\square C_\square$.

The terms in \eqref{eq:global-parametrix-inverse} have a local
description. A walk of length $n$ is a sequence
$\omega=(\square_0,\ldots,\square_n)$ with successive centers at
distance at most $M$. Its term is
\[
 S_\omega=(h_0S_0h_0)(C_1S_1h_1)\cdots(C_nS_nh_n),
\]
where the subscript $i$ denotes the cube $\square_i$.
Other sequences give zero because of the support of the commutator.
Its connection variables lie in
$\bigcup_{i=0}^nW_{\square_i}$.

\begin{proposition}\label{prop:walk-propagator-bounds}
For $p\ge1$ with $q(p)=54\pi C_Va'p/M\le1/2$,
\begin{align}
 \sup_z\sum_{z'}\exp(\gamma d(z,z'))
       \sum_{|\omega|=n}p^{n+1}
           \|\mathbf1_{\Delta_z}S_\omega\mathbf1_{\Delta_{z'}}\|
       &\le16a'pw_k(y)q(p)^n,\label{eq:walk-row-bound}\\
 \sup_z\sum_{z'}\exp(\gamma d(z,z'))
       \sum_{|\omega|=n}p^{n+1}
           \|\mathbf1_{\Delta_z}\delta_{\alpha,A}S_\omega
                            \mathbf1_{\Delta_{z'}}\|
       &\le152a'pq(p)^n.\label{eq:walk-holder-bound}
\end{align}
The norms are sup operator norms. At $k=0$ the first bound holds
also with the supremum of the normalized kernel over the two clusters.
In particular, $S_y(A)$ is analytic on $\widetilde{\mathcal G}_k$ and
\begin{equation}\label{eq:global-propagator-bounds}
 N_\gamma(S_y(A))\le32a'w_k(y),\qquad
 N_\gamma(\delta_{\alpha,A}S_y(A))\le304a'.
\end{equation}
These constants are independent of the level, the side, the running
mass in $[0,1]$, and the connection in the stated domain.
\end{proposition}

\begin{proof}
The first factor in a walk has norm at most $2a'w_k(y)$.
Every additional factor has norm at most $2\pi C_Va'/M$ by
\eqref{eq:local-inverse-bounds} and \eqref{eq:cutoff-commutator-bound}.
For a fixed output block, at most eight initial cutoffs contribute;
there are at most $27$ choices at each further step. Apply
submultiplicativity row by row and sum these $8\cdot27^n$ bounds.
This proves \eqref{eq:walk-row-bound}.

For the covariant quotient, only the factor at the left end of the
product is differentiated. The product identity is
\[
 \delta_{\alpha,A}(h_\square v)(x,x')
 =h_\square(c)\delta_{\alpha,A}v(x,x')
   +d_X^{-\alpha}(h_\square(c')-h_\square(c))U_A(x,x')v(x').
\]
The first contribution has localized norm at most $18a'$.
The second has coefficient at most $3\pi\exp(1)/M$, and evaluation
at the second point costs at most $27\exp(\gamma)$. With the stated
choice of $M$ this contribution is at most $a'$. Thus
$N_\gamma(\delta_{\alpha,A}h_\square S_\square)\le19a'$.
For intracluster pairs the second term vanishes, since $c=c'$.
There are again at most eight possible initial cubes, including when
the second point of a short pair enters the cutoff support. The count
$8\cdot27^n$ now gives \eqref{eq:walk-holder-bound}.

At level zero keep the kernel norm in the rightmost local inverse;
the other factors are controlled in operator norm. The same
submultiplicative row estimate proves the kernel assertion.
Summing the geometric bounds with $p=1$ gives
\eqref{eq:global-propagator-bounds}. For a fixed torus the sum of the
operator norms over all walks is finite; its tail of lengths at least
$j$ is bounded by $4a'(\ell/M)^3 2^{-j}$. Each walk term is analytic
in the finitely many connection variables it reads. The uniform tail
bound therefore proves analyticity of the sum. The local bounds in
\eqref{eq:walk-row-bound} contain no factor depending on $\ell$.
\end{proof}


\begin{proof}[Proof of \Cref{thm:chamber-locality}]
Choose $\gamma,a'$ as in \eqref{eq:zero-row-bounds}.
The constants in the free estimates depend only on $\rho,L,\alpha$.
The two charge conditions give \eqref{eq:local-inverse-bounds},
and $q(1)\le1/2$ gives the convergent parametrix.
\Cref{prop:walk-propagator-bounds} supplies the analytic inverse and
both asserted row estimates, including the normalized kernel estimate
at level zero.
\end{proof}
\section{Localized propagators and coarse covariances}\label{sec:fermion-primitives}

We use the local integration construction of \cite[Section~11]{QED3Manuscript}. The walk expansion provides its localization. A parameter is attached to each $M$-cube traversed
by a walk, together with its one-layer neighborhood. Setting parameters
to zero then separates the operator into components. The coarse
covariance is obtained from this expansion by an exact resolvent
identity. At parameter value one it agrees with the exact covariance;
its local terms expand the determinant and the interaction.
We work at mesh $\eta_k=L^{-k}$ on the torus of side $\ell$; this is
the side denoted by $R$ in the local inverse construction.
Throughout this section the mass satisfies $0<\mu\le1$,
the range of the renormalization data.
Transposition is with respect to the bilinear lattice pairings:
for $O:E\to E'$, $\langle v,Ou\rangle_{E'}=
\langle O^Tv,u\rangle_E$. These pairings use the counting,
chamber, or coarse lattice measure of the respective space.
No complex conjugation enters this operation.

\subsection{Weakened walks and polymer propagators}\label{subsec:weakened-walks}

The resolvent parameter ranges over $Y_0=\R$ at level zero and
$Y_k=\R\cup\{\infty\}$ at $k\ge1$, with $-\infty=\infty$
as in the preceding resolvent construction.

Let $\pi_k$ be the partition into $M$-cubes. For integers $j\ge0$, write
$\square^{\sim j}$ for the union of partition cubes whose centers
are within periodic sup distance $jM$ of the center of $\square$.
When this enlargement does not wrap, it has side $(2j+1)M$.
Let $\mathcal P_k$ be the set of
nonempty unions of these cubes connected through faces. For $X\in\mathcal P_k$
write $|X|_M$ for its number of cubes. With the net and length
conventions of \cref{app:polymer-geometry}, set
\[
 M d_M(X)=\inf\{\operatorname{length}(\mathcal N):
       \mathcal N\text{ is a net meeting every closed cube of }X\}.
\]
For a face-connected union of cubes, joining successive cube centers gives
$d_M(X)\le |X|_M-1$.

For a walk $\omega=(\square_0,\ldots,\square_n)$ from
\cref{subsec:propagator-walks}, put
\[
 X_\omega=\bigcup_{i=0}^n\square_i^{\sim1},\qquad
 s_\omega=\prod_{\square\subset X_\omega}s_\square,
 \qquad |s_\square|\le M^{\alpha_0},\quad \alpha_0=1/54.
\]
The union $X_\omega$ is a polymer and
\begin{equation}\label{eq:walk-polymer-size}
 |X_\omega|_M\le27(n+1),\qquad d_M(X_\omega)\le27(n+1).
\end{equation}
Indeed, each one-layer neighborhood contains $27$ cubes. Successive
centers differ by at most $M$ in each coordinate, so their neighborhoods
are joined through faces. The preceding bound on $d_M$ applies.
Every walk, including a walk of length zero, carries the parameters of
all cubes in $X_\omega$.

The left parametrix gives terms
\[
 S'_{y,\omega}=(h_nS_n)(-C_nh_{n-1}S_{n-1})\cdots
                 (-C_1h_0S_0)h_0.
\]
Here, as for $S_{y,\omega}$, $S_i$ is the local inverse on
$\square_i$ and $C_i=-[F_y,h_i]$. Define
\begin{align}
 S_y(s,A)&=\sum_\omega s_\omega S_{y,\omega}(A),&
 S_y(X,A)&=\sum_{X_\omega=X}S_{y,\omega}(A),\label{eq:weakened-walk-definitions}\\
 S'_y(s,A)&=\sum_\omega s_\omega S'_{y,\omega}(A),&
 S'_y(X,A)&=\sum_{X_\omega=X}S'_{y,\omega}(A).\notag
\end{align}
Their averages are denoted by $S_y^{\mathrm{sy}}$. The superscript
$\#$ below means any of these three choices. Put
$\kappa'=\kappa_{\mathrm w}+\kappa_0$, where
$\kappa_0>0$ and $\kappa_{\mathrm w}=\kappa+\kappa_0+3$, with $\kappa>0$, are fixed polymer decay rates.
Choose $M$ large enough that
\begin{equation}\label{eq:weakened-walk-contraction}
 q(M^{27\alpha_0})\le\tfrac12,\qquad
 q(e^{27\kappa'})\le\tfrac12,
 \qquad q(p)=54\pi C_Va'p/M.
\end{equation}

\begin{proposition}\label{prop:weakened-polymer-propagators}
For $A\in\widetilde{\mathcal G}_k$, $y\in Y_k$, and
$s\in\prod_{\pi_k}\{|s_\square|\le M^{\alpha_0}\}$, the series
\eqref{eq:weakened-walk-definitions} converge absolutely in $N_\gamma$,
uniformly in $(s,A)$, and
\begin{align}
 N_\gamma(S_y^\#(s,A))&\le32a'M^{27\alpha_0}w_k(y),
                                       \label{eq:weakened-propagator-bound}\\
 N_\gamma(\delta_{\alpha,A}S_y^\#(s,A))&\le304a'M^{27\alpha_0},\notag\\
 \|\mathbf1_{\Delta_z}S_y^\#(X,A)\mathbf1_{\Delta_{z'}}\|
   &\le32a'e^{27\kappa'}w_k(y)
       e^{-\gamma d(z,z')-\kappa'd_M(X)}.
                                       \label{eq:polymer-propagator-bound}
\end{align}
The norm in the last line is the sup operator norm. At $k=0$ it can
also be replaced by the supremum of the normalized kernel over the two
clusters, with $w_0(y)=(1+|y|)^{-1}$.
At parameter value one,
\begin{equation}\label{eq:polymer-propagator-sum}
 S_y^\#(1,A)=S_y(A)=\sum_{X\in\mathcal P_k}S_y^\#(X,A).
\end{equation}
The weakened operators are continuous on the closed parameter
polydisc, jointly analytic in its interior and in $A$, and analytic in
$A$ for each fixed parameter value, including boundary values.
For finite real $y$ the families are continuous in $y$.

The operator $S_y^\#(X,A)$ is supported on $X$ in both arguments and
depends on $A$ only on bonds with endpoints in $X$. If parameters vanish
off a union $Y$ of $M$-cubes, then each component of $Y$ is preserved by
$S_y^\#(s,A)$. On such a component the operator depends only on its
parameters and its connection. Its row at $x=(c,g)$ vanishes unless all
$M$-cubes meeting $K_c=\{c':|c'-c|_\infty\le2L\}$ have active
parameters; the corresponding assertion holds for columns.
\end{proposition}

\begin{proof}
By \eqref{eq:walk-polymer-size}, $|s_\omega|\le p_s^{n+1}$ with
$p_s=M^{27\alpha_0}$. Apply the walk row estimate with this weight.
The length-$n$ contribution is bounded by
$16a'p_sw_k(y)2^{-n}$; its H\"older row estimate is bounded by
$152a'p_s2^{-n}$. Summing these bounds proves the first two estimates.
For a fixed polymer use
\[
 \mathbf1[X_\omega=X]\le
       e^{-\kappa'd_M(X)}p_X^{n+1},\qquad p_X=e^{27\kappa'}.
\]
The same row estimate with $p_X$ proves
\eqref{eq:polymer-propagator-bound}. The left walk estimates follow
from the left parametrix by the same weighted product estimate; taking
the average preserves all bounds.

On a fixed torus, summing the row bounds over its finitely many unit
blocks also bounds the sum of the $N_\gamma$ norms of the individual
terms. The tails decay geometrically, uniformly in $(s,A)$. The series
may therefore be regrouped by $X_\omega$. At $s=1$ they are the right
and left inverse expansions, so both equal $S_y(A)$. This proves
\eqref{eq:polymer-propagator-sum}. Each finite sum is analytic in its
connection variables and polynomial in $s$. The uniform tail estimates
give the stated analyticity and continuity.
The local inverse is continuous at every real resolvent parameter.
Its walk terms inherit this property, and their tails are uniform in
$y$ since $w_k(y)\le1$. The sums are therefore continuous in real $y$.

A walk term is a product of local factors on the neighborhoods
$\square_i^{\sim1}$. Its row and column cutoffs, its commutators, and
all connection variables consequently lie in $X_\omega$. If the first
row is nonzero, its photon site satisfies
$|c-z_{\square_0}|_\infty<3M/4$. Since $3M/4+2L<3M/2$, every
$M$-cube meeting $K_c$ is in $\square_0^{\sim1}$, and its parameter
occurs in $s_\omega$. The last cutoff gives the column assertion;
the order is reversed for a left walk. Finally, a surviving walk has
every cube of $X_\omega$ in $Y$. Since $X_\omega$ is face-connected,
it lies in a single component of $Y$. Its local factors use only the
connection in that component. This proves all the locality statements.
\end{proof}
The lattice symmetries used in this section have coordinate sign
changes and translations in $M\Z^3$. Their spin lifts lie in the
group generated by $\gamma_1,\gamma_2,\gamma_3,\gamma_5$; they act
on $(c,g)$ by $(rc+a,rg)$, where
$S\gamma_\mu S^{-1}=r_{\mu\mu}\gamma_\mu$ determines $r$.
Write $\widetilde g=(S,a)$ for such a symmetry.
The step group restricts $a$ to
$LM\Z^3$. A lift commuting with $\gamma_5$ preserves $y$; a lift
anticommuting with it sends $y$ to $-y$ at nonneutral values.
Gauge covariance follows term by term: the conjugating phases of
adjacent factors cancel. A lattice symmetry permutes cubes and walks,
and hence acts on the parameters by
$(\widetilde g s)_\square=s_{\widetilde g^{-1}\square}$.
Transposition reverses the order of
the local factors and interchanges the right and left expansions.
More precisely, with $\mathsf C=i\gamma_2$ one has
$F_y(A)^T=\mathsf C^{-1}F_{-y}(-A)\mathsf C$. Therefore
\begin{equation}\label{eq:weakened-transpose}
 S_y(s,A)^T=\mathsf C^{-1}S'_{-y}(s,-A)\mathsf C.
\end{equation}
Thus at the neutral values $y=0,\infty$ for $k\ge1$, and $y=0$
for $k=0$, the symmetrized family has the same transpose covariance as
the exact inverse. This is the reason for using it in the coarse
covariance and determinant.

\subsection{The fine and coarse resolvents}\label{subsec:fine-coarse-resolvent}

For $k\ge1$ let
\[
 \mathcal Q_k=Q_k(A)J^*,\qquad
 \mathcal Q_k^\sharp=JQ_k^T(-A),\qquad
 P_L=Q^T(-A)Q(A).
\]
Here $Q_k$ averages the fine lattice onto the unit lattice, and $Q$
averages the unit lattice onto the lattice of spacing $L$.
The notation $\sharp$ denotes the displayed algebraic transpose with
reversed connection; for complex $A$ it is not a Hermitian adjoint.
Here $B(z')=E_{z'}\cap\Z^3
=\{z\in\Z^3:|z-z'|_\infty<L/2\}$ is the unit-lattice
part of the $L$-block in \cref{subsec:blocked-zero-inverse};
the blocks are taken modulo the torus periods.
To specify the phases, if $c\in\Delta_z$ and $z\in B(z')$, set
\[
 \varphi_A(z,c)=e^{ie_k\eta_k(\tau_kA)(z,c)},\qquad
 U_A(z',z)=e^{ie_k\eta_k(\tau A)(z',z)}.
\]
The path average $\tau A$ is the arithmetic mean of the oriented bond
sums in the six coordinate orders, and $\tau_kA$ is their sum along
the nested block centers. Then
\begin{align*}
 (\mathcal Q_kf)(z)&=\eta_k^3\sum_{c\in\Delta_z}
                \varphi_A(z,c)\frac1{48}\sum_gf(c,g),\\
 (\mathcal Q_k^\sharp v)(c,g)&=\varphi_A(z,c)^{-1}v(z),\\
 (P_Lv)(z)&=U_A(z',z)^{-1}L^{-3}
                    \sum_{z''\in B(z')}U_A(z',z'')v(z'').
\end{align*}
The nesting of the odd blocks and cancellation of the phases give
\begin{equation}\label{eq:block-algebra}
 \mathcal Q_k\mathcal Q_k^\sharp=1,\quad P_L^2=P_L,\quad
 \mathcal Q_k^\sharp\mathcal Q_k=\Pi_1(A),\quad
 \mathcal Q_k^\sharp P_L\mathcal Q_k=\Pi_L(A).
\end{equation}
All these operators commute with constant spin matrices.

Write $S_k=S_\infty(A)$ and define maps on the unit spinors by
\begin{align}
 H_k&=b_kS_k\mathcal Q_k^\sharp,&
 \mathfrak H_k&=b_k\mathcal Q_kS_k,\notag\\
 D_k&=b_k-b_k^2\mathcal Q_kS_k\mathcal Q_k^\sharp,&
 M_k&=b_k-D_k.\label{eq:coarse-dirac-definition}
\end{align}
For real $y$ put
\[
 \mathcal B_{k,y}=(b_k+i\gamma_5y)^{-1}(1-P_L)
                   +(b_k+b_L+i\gamma_5y)^{-1}P_L.
\]
These $4\times4$ inverses exist since
$(t+i\gamma_5y)^{-1}=(t-i\gamma_5y)/(t^2+y^2)$ for $t>0$.
The idempotent decomposition and \eqref{eq:block-algebra} give
\begin{equation}\label{eq:fine-resolvent-perturbation}
 \mathcal B_{k,y}^{-1}=b_k+b_LP_L+i\gamma_5y,\qquad
 F_y=F_\infty-b_k^2\mathcal Q_k^\sharp
                       \mathcal B_{k,y}\mathcal Q_k.
\end{equation}

\begin{proposition}\label{prop:fine-coarse-resolvents}
For $k\ge1$, $A\in\widetilde{\mathcal G}_k$, and real $y$, the
operator $D_k+b_LP_L+i\gamma_5y$ is invertible and its inverse is
\begin{equation}\label{eq:coarse-resolvent-formula}
 \Gamma_{k,y}=\mathcal B_{k,y}
       +b_k^2\mathcal B_{k,y}\mathcal Q_kS_y
                         \mathcal Q_k^\sharp\mathcal B_{k,y}.
\end{equation}
Moreover,
\begin{equation}\label{eq:resolved-fine-covariance}
 S_y=S_k+H_k\Gamma_{k,y}\mathfrak H_k.
\end{equation}
In particular, $\Gamma_k=\Gamma_{k,0}$ exists and
$\delta Z_k(A)=\det(D_k+b_LP_L)\ne0$.
\end{proposition}

\begin{proof}
Set $Z=b_k^2\mathcal Q_k^\sharp\mathcal B_{k,y}\mathcal Q_k$.
The two fine inverses are constructed in
\cref{subsec:propagator-walks} and recovered at parameter value
one in \cref{prop:weakened-polymer-propagators}. Equation
\eqref{eq:fine-resolvent-perturbation} gives both identities
\[
 S_y-S_k=S_kZS_y=S_yZS_k.
\]
Since $D_k+b_LP_L+i\gamma_5y
=\mathcal B_{k,y}^{-1}-b_k^2\mathcal Q_kS_k\mathcal Q_k^\sharp$,
multiplication by the right side of
\eqref{eq:coarse-resolvent-formula} gives the identity: the additional
term is
$b_k^2\mathcal Q_k(S_y-S_k-S_kZS_y)
\mathcal Q_k^\sharp\mathcal B_{k,y}=0$.
The product in the other order uses $S_y-S_k-S_yZS_k=0$.
This proves the inverse formula without assuming that the coarse
operator is invertible. Substituting the formula in
$H_k\Gamma_{k,y}\mathfrak H_k$ gives
\[
 S_kZS_k+S_kZS_yZS_k
     =S_kZ(S_k+S_yZS_k)=S_kZS_y=S_y-S_k.
\]
At $y=0$ the coarse operator is an invertible map on the finite
dimensional unit spinor space, so its determinant is nonzero.
\end{proof}

At level zero this algebra is replaced by a single chamber
integration. The operators $b_0$, $S_\infty$, $\mathcal B_{0,y}$, and
$M_0$ are not defined. Instead,
\begin{equation}\label{eq:level-zero-covariance}
 F_0=D_A+\mu+b_LJP_LJ^*,\qquad
 \Gamma_0=F_0^{-1},\qquad \delta Z_0(A)=\det F_0\ne0.
\end{equation}
Here $\Gamma_0$ acts on $192$ components per photon site, whereas
$\Gamma_k$ for $k\ge1$ acts on four components per unit site.

\subsection{Localized coarse covariances}\label{subsec:localized-coarse-covariance}

The weakened families for $k\ge1$ are
\begin{align}
 H_k(s,A)&=b_kS_\infty(s,A)\mathcal Q_k^\sharp,
                                          \label{eq:local-coarse-families}\\
 \Gamma_k(s,A)&=\mathcal B_{k,0}
    +b_k^2\mathcal B_{k,0}\mathcal Q_kS_0^{\mathrm{sy}}(s,A)
                 \mathcal Q_k^\sharp\mathcal B_{k,0},\notag\\
 M_k(X,A)&=b_k^2\mathcal Q_kS_\infty^{\mathrm{sy}}(X,A)
                  \mathcal Q_k^\sharp.\notag
\end{align}
The unprimed polymer propagators $S_{k,y}(X,A)=S_y(X,A)$ are used
in the determinant integral. At level zero let
$\operatorname{Cub}(c)$ be the $M$-cubes having a lattice site in
common with $K_c$, and set
\begin{equation}\label{eq:local-level-zero-families}
 \begin{aligned}
 (H_0(s,A)u)(c,g)&=
       \left(\prod_{\square\in\operatorname{Cub}(c)}s_\square\right)u(c,g),\\
 \Gamma_0(s,A)&=S_0^{\mathrm{sy}}(s,A),\qquad
 S_{0,y}(X,A)=S_y(X,A).
 \end{aligned}
\end{equation}
At $s=1$ these are the exact operators; also
$\sum_XM_k(X,A)=M_k(A)$. The exponential formulas for the averaging
phases and \cref{prop:weakened-polymer-propagators} give
analyticity on the same domains.

For use in the estimates, let
$\mathfrak i(A)=e_k\|\operatorname{Im}A\|_\infty$.
The small-connection domain and the charge bound give
$\mathfrak i(A)\le1/(6L)$; reduce $\gamma$, if necessary, so that
$\gamma L\le1/6$. The path lengths imply
\begin{equation}\label{eq:coarse-phase-bounds}
 |\varphi_A(z,c)^{\pm1}|\le e^{3\mathfrak i(A)/2},\qquad
 \|\mathcal Q_k\|_{\infty\to\infty},
 \|\mathcal Q_k^\sharp\|_{\infty\to\infty}
       \le e^{3\mathfrak i(A)/2}.
\end{equation}
The sup norms here are formed with the Euclidean spinor norm.
Each unit block has chamber measure one. In particular, for a fine
operator $O$,
\begin{equation}\label{eq:block-compression-bound}
 |(\mathcal Q_kO\mathcal Q_k^\sharp)(z,z')|_2
       \le \exp(1)\,\|\mathbf1_{\Delta_z}O\mathbf1_{\Delta_{z'}}\|.
\end{equation}
Indeed, apply the two averaging bounds to a unit spinor supported at
$z'$; their product is $e^{3\mathfrak i(A)}\le \exp(1)$.

\begin{proposition}\label{prop:coarse-local-family-bounds}
The families \eqref{eq:local-coarse-families} satisfy
\begin{align}
 \|H_k(s,A)\|_{\infty\to\infty}
       &\le32\exp(1)b_ka'M^{27\alpha_0},\label{eq:weakened-H-bound}\\
 \|\delta_{\alpha,A}H_k(s,A)\|_{\infty\to\infty}
       &\le304\exp(1)b_ka'M^{27\alpha_0},\notag\\
 |\Gamma_k(s,A)(z,z')|_2
       &\le C_\Gamma M^{27\alpha_0}e^{-\gamma d(z,z')},
                                             \label{eq:weakened-Gamma-bound}\\
 |M_k(X,A)(z,z')|_2
       &\le32\exp(1)b_k^2a'e^{27\kappa'}
                     e^{-\gamma d(z,z')-\kappa'd_M(X)},\notag
\end{align}
where
\[
 C_\Gamma=\frac{2(1+2\exp(1))}{b_-}+128\exp(1)(1+2\exp(1))^2a'.
\]
The H\"older estimate is on short pairs, including the distinct
sites of one cluster. The families $H_k(s,A)$ and $\Gamma_k(s,A)$
preserve every component of the active $M$-cubes and depend there only
on its parameters and connection. The row of $H_k(s,A)$ at $(c,g)$
vanishes unless all cubes meeting $K_c$ are active. The polymer
$M_k(X,A)$ is supported on the unit sites of $X$ and depends only on
the connection in $X$.

At level zero the normalized kernel of $\Gamma_0(s,A)$ is bounded by
$32a'M^{27\alpha_0}e^{-\gamma d(c,c')}$, and
\eqref{eq:polymer-propagator-bound} bounds $S_{0,y}(X,A)$ with its
factor $(1+|y|)^{-1}$. The multiplier $H_0(s,A)$ has sup norm
at most $M^{8\alpha_0}$ and equals the identity at $s=1$.
Its row at $(c,g)$ vanishes unless every cube meeting $K_c$
is active. There are no short pairs of distance less than one
at this level.
\end{proposition}

\begin{proof}
The first two estimates follow by composing the weakened fine bounds
with \eqref{eq:coarse-phase-bounds}. The third and fourth use
\eqref{eq:block-compression-bound}. For the covariance it remains to
control the two factors $\mathcal B_{k,0}$. Its kernel is
\[
 \mathcal B_{k,0}(z,z')=b_k^{-1}\mathbf1[z=z']1_4
      -\frac{b_L}{b_k(b_k+b_L)}P_L(z,z').
\]
It vanishes unless $z,z'$ belong to one $L$-block. Since
$|P_L(z,z')|_2\le\exp(1)L^{-3}$ on that block, its row and column sums
are at most $(1+2\exp(1))/b_k$. In a nonzero term of
$\mathcal B_{k,0}G\mathcal B_{k,0}$,
where $G=\mathcal Q_kS_0^{\mathrm{sy}}(s,A)\mathcal Q_k^\sharp$,
the two intermediate sites are within $L$ of the exterior sites.
Therefore
\[
 e^{-\gamma d(z_1,z_2)}\le
      e^{2\gamma L}e^{-\gamma d(z,z')}
      \le4e^{-\gamma d(z,z')}.
\]
The fine bound and \eqref{eq:block-compression-bound} give
$|G(z_1,z_2)|_2\le32\exp(1)a'M^{27\alpha_0}e^{-\gamma d(z_1,z_2)}$.
Multiplying the row and column sums cancels the factor $b_k^2$ and
gives $128\exp(1)(1+2\exp(1))^2a'M^{27\alpha_0}$.
The local term $\mathcal B_{k,0}$ is bounded by
$2(1+2\exp(1))b_-^{-1}e^{-\gamma d(z,z')}$, proving
\eqref{eq:weakened-Gamma-bound}.

The odd block partitions are nested, and every unit block and every
$L$-block belongs to a single $M$-cube. Moreover, the paths entering
$\varphi_A(z,c)$ lie in $\Delta_z$, and the paths entering
$U_A(z',z)$ lie in the fine block of side $L$ centered at $z'$:
each coordinate of a rectilinear path remains between those of its
endpoints. Consequently the averaging maps and $\mathcal B_{k,0}$
preserve every union of $M$-cubes and depend only on its connection.
Composing them with the component-local fine families proves the
locality assertions. This applies also to the parameter-independent
term $\mathcal B_{k,0}$ in $\Gamma_k(s,A)$.
For the row of $H_k$, the terminal average does not change the first
row index of $S_\infty(s,A)$, so the stronger vanishing statement is
preserved. At level zero the assertions for $\Gamma_0$ follow
from the fine kernel bound and locality with $\eta_0=1$.

For $H_0$, lift $c$ and $K_c$ to the lattice. In each coordinate
$K_c$ is an interval of $4L+1$ sites. Since $M\ge L^2$ and $L\ge5$,
this interval meets at most two consecutive $M$-blocks. The cubes
meeting $K_c$ are the product of these three sets of blocks.
They number at most eight, form a face-connected set, and include
the cube of $c$. Thus the defining product is a monomial of degree
at most eight, bounded by $M^{8\alpha_0}$. If one of its cubes is
inactive, its factor is zero. If all are active, their connectedness
puts them in the active component containing $c$; hence changing
parameters outside that component has no effect on the row or
column at $(c,g)$. The multiplier is scalar in spin and chamber
labels and independent of the connection. These observations give
its analyticity, gauge covariance, and component locality.
\end{proof}

These families have the gauge and lattice covariance of their exact
counterparts. For $\Gamma_k$ and $M_k$, transposition follows from
\eqref{eq:weakened-transpose} at a neutral value and the scalar
averaging phases. The barred field map is
$\bar H_k(s,A)=\mathsf C^{-1}H_k(s,-A)\mathsf C$.
Its bounds follow from those for $H_k$ because $\mathsf C$ is unitary
and $\mathsf C\gamma_\mu\mathsf C^{-1}=-\gamma_\mu^T$.
When component sup norms are used in the Grassmann substitution,
the Euclidean operator bounds above acquire at most the factor two
from $|v|_c\le|v|_2\le2|v|_c$.

\section{Weighted Grassmann block integration}\label{sec:weighted-integration}

\subsection{Weighted Grassmann integration}\label{subsec:weighted-grassmann}

Let $\phi_i(t)$, $1\le i\le p$, be independent Grassmann families
on finite ordered index sets. If $\phi_i(I_i)$ is the ordered product
over $I_i$, write
\[
 F=\sum_{I_1,\ldots,I_p}c_F(I_1,\ldots,I_p)
                 \phi_1(I_1)\cdots\phi_p(I_p),\qquad
 \|F\|_{\mathsf w}=\sum_I|c_F(I)|\prod_i\mathsf w_i^{|I_i|}.
\]
The positive measures $\nu_i(t)$ on the index sets enter the kernels,
not this coefficient norm. Thus an antisymmetric kernel expansion
\[
 F=\sum_n\prod_i\frac1{n_i!}\int\dd t\,K_n(t)
                       \prod_i\prod_{j=1}^{n_i}\phi_i(t_j^i)
\]
has norm
\begin{equation}\label{eq:grassmann-kernel-norm}
 \|F\|_{\mathsf w}
     =\sum_n\prod_i\frac{\mathsf w_i^{n_i}}{n_i!}
                      \int\dd t\,|K_n(t)|.
\end{equation}
For arbitrary kernels the right side is an upper bound. Products of
ordered monomials vanish on repeated generators and otherwise multiply
their weights, so
\begin{equation}\label{eq:grassmann-product-bound}
 \|FG\|_{\mathsf w}\le\|F\|_{\mathsf w}\|G\|_{\mathsf w},
 \qquad \|e^F-1\|_{\mathsf w}\le e^{\|F\|_{\mathsf w}}-1
       \quad(F\text{ even}).
\end{equation}

For an additional finite nonempty mode set $\mathcal I$ with positive weights
$\nu(t)$, introduce $W(t),\bar W(t)$ and set
\[
 \langle\bar W,OW\rangle_\nu
    =\sum_{t,t'}\nu(t)\nu(t')\bar W(t)O(t,t')W(t'),
 \qquad \mathrm{Top}=\prod_{t\in\mathcal I}W(t)\bar W(t).
\]
The kernel represents $O$ with measure $\nu$; its matrix entry is
$O(t,t')\nu(t')$. Define $\int\mathrm DW$ as extraction of the
coefficient of $\mathrm{Top}$, multiplied by
$\prod_t\nu(t)^{-1}$, and define
\begin{align}
 \int\dd\mu_I(W)F
     &=\int\mathrm DW\,e^{-\langle\bar W,W\rangle_\nu}F,
                                            \label{eq:weighted-Berezin}\\
 \int\dd\mu_C(W)F
     &=(\det C)\int\mathrm DW\,
               e^{-\langle\bar W,C^{-1}W\rangle_\nu}F
                    \quad(C\text{ invertible}).\notag
\end{align}
These are algebra-valued functionals when other Grassmann generators
are present. The order of modes in $\mathrm{Top}$ is immaterial because
its factors are even.

\begin{proposition}\label{prop:weighted-Berezin-bounds}
Put $w_W=\max_t\nu(t)^{-1/2}$. Then
\begin{equation}\label{eq:weighted-integration-bound}
 \left\|\int\dd\mu_I(W)F\right\|_{\mathsf w}
                  \le\|F\|_{(\mathsf w,w_W)}.
\end{equation}
The integral factors for polynomials involving disjoint sets of modes.
For an invertible covariance,
\begin{align}
 \int\dd\mu_C(W)F(\bar W,W)
     &=\int\dd\mu_I(W)F(\bar W,CW),\label{eq:covariance-substitution}\\
 \int\dd\mu_C(W)W(t)\bar W(t')&=C(t,t'),\notag\\
 \int\mathrm DW\,e^{-\langle\bar W,OW\rangle_\nu}
     &=\det O\quad\text{for every }O.\label{eq:weighted-determinant}
\end{align}
Finally, consider a linear substitution
$\phi_i\mapsto\sum_jO_{ij}\theta_j$, and put
\[
 \|O_{ij}\|_{1,\infty}
     =\sup_t\sum_s|O_{ij}(t,s)|\nu'_j(s).
\]
If $\sum_j\|O_{ij}\|_{1,\infty}\mathsf w'_j\le\mathsf w_i$ for
each $i$, then its image $F'$ obeys
$\|F'\|_{\mathsf w'}\le\|F\|_{\mathsf w}$.
Only rows corresponding to generators occurring in $F$ need satisfy
this bound.
\end{proposition}

\begin{proof}
The identity density factors as
\[
 e^{-\langle\bar W,W\rangle_\nu}
            =\prod_t(1+\nu(t)W(t)\bar W(t)).
\]
An ordered monomial $W(J)$ contributes to its integral only if $J$
contains both generators of every mode it meets. In that case its
integral is its reordering sign times
$\prod_{(t,0)\in J}\nu(t)^{-1}$; otherwise it is zero. Its modulus
is therefore at most $\prod_{u\in J}\nu(u)^{-1/2}\le w_W^{|J|}$.
Apply this coefficientwise to obtain
\eqref{eq:weighted-integration-bound}. For disjoint mode sets a
nonzero integral requires each monomial separately to be paired.
Their degrees in $W$ are then even, so the signs from exchanging
external generators cancel. The paired-mode formula gives
factorization. It also shows that modes absent from the integrand may
be removed.

A linear change $W\mapsto OW$ multiplies the coefficient of the top
monomial by $\det O$, since the matrix of the change is
$O(t,t')\nu(t')$. Applied to the identity density this proves
\eqref{eq:weighted-determinant}, without an invertibility assumption.
For invertible $C$, use $W\mapsto CW$, leaving $\bar W$ fixed, in
the density with covariance $C$. This proves
\eqref{eq:covariance-substitution} and normalization. Its two-point
function is
\[
 \sum_s C(t,s)\nu(s)\int\dd\mu_I(W)W(s)\bar W(t')
       =\sum_s C(t,s)\nu(s)\frac{\mathbf1[s=t']}{\nu(t')}
       =C(t,t').
\]
Thus it is the kernel, rather than the matrix entry.

Under the stated substitution, the norm of the image of
$\phi_i(t)$ is at most
\[
 \sum_j\mathsf w'_j\sum_s|O_{ij}(t,s)|\nu'_j(s)\le\mathsf w_i
\]
by hypothesis. Submultiplicativity bounds
the image of each monomial by its original weight; summing its
coefficients proves the last assertion, with the same proof on the
subalgebra generated by the rows that occur.
\end{proof}

For linear combinations $v_1,\ldots,v_{2r}$ of the $W,\bar W$,
Wick's formula gives
\begin{equation}\label{eq:weighted-wick}
 \int\dd\mu_I(W)v_1\cdots v_{2r}
      =\operatorname{Pf}\left(\int\dd\mu_I(W)v_iv_j\right)_{i,j}.
\end{equation}
Odd products integrate to zero. To fix the signs, let $\partial_u$
be the left Grassmann derivative. It satisfies
$\partial_u(vG)=(\partial_uv)G-v\partial_uG$, and its Berezin
integral is zero. Differentiating the factored Gaussian density gives
\begin{align*}
 \int\dd\mu_I(W)W(t)G
      &=\nu(t)^{-1}\int\dd\mu_I(W)\partial_{\bar W(t)}G,\\
 \int\dd\mu_I(W)\bar W(t)G
      &=-\nu(t)^{-1}\int\dd\mu_I(W)\partial_{W(t)}G.
\end{align*}
Consequently, if $I(v)=\int\dd\mu_I(W)v$,
\[
 I(v_1\cdots v_q)=\sum_{j=2}^q(-1)^jI(v_1v_j)
                         I(v_2\cdots\widehat v_j\cdots v_q).
\]
This is the expansion of the Pfaffian along its first row and proves
\eqref{eq:weighted-wick} by induction. Formula
\eqref{eq:covariance-substitution} gives the corresponding Wick rule
for $\mu_C$.

In the renormalization step the modes have weight one for $k\ge1$.
At $k=0$ they are $(c,g,\beta)$ with weight $1/48$. Thus the
integration weight is
\begin{equation}\label{eq:first-step-Grassmann-weight}
 w_W=1\quad(k\ge1),\qquad w_W=\sqrt{48}\quad(k=0).
\end{equation}
In particular the first-step identity covariance has two-point
function $48\mathbf1[(c,g,\beta)=(c',g',\beta')]$.
The pair index used for a H\"older field has measure
$(\eta_k^3/48)^2$. In a substitution involving
$\delta_{\alpha,A}H_k$, its scalar component row sum is computed
with the input measure, as in
\cref{prop:weighted-Berezin-bounds}. Internal-cluster
differences are therefore controlled by the H\"older bound with the
same field weight as other pairs.

\subsection{The Gaussian block identity}\label{subsec:Gaussian-block-identity}

For this identity use a common notation for the integration space:
at $k\ge1$ let $\mathcal D_k=D_k$ and
$\mathcal A_k(A)=Q(A)$; at $k=0$ let
$\mathcal D_0=D_A+\mu$ and $\mathcal A_0(A)=Q(A)J^*$.
The integration measure is respectively counting measure or the
chamber measure $1/48$. The target is the four-component spinor
space on the lattice of spacing $L$, with measure $L^3$.
Set
\[
 \mathcal A_k^\sharp(A)=\mathcal A_k(-A)^T,\qquad
 K_k=\mathcal D_k+b_L\mathcal A_k^\sharp\mathcal A_k,
 \qquad\Gamma_k=K_k^{-1}.
\]
Thus $K_k=D_k+b_LP_L$ for $k\ge1$ and $K_0=F_0$.
Define the quadratic operator on the coarse lattice by
\[
 D_{k+1}^0=b_L-b_L^2\mathcal A_k\Gamma_k
                                       \mathcal A_k^\sharp.
\]
Let $\Phi,\bar\Phi$ be external coarse Grassmann fields and put
\[
 \Psi^c=b_L\Gamma_k\mathcal A_k^\sharp\Phi,\qquad
 \bar\Psi^c=b_L\Gamma_k^T\mathcal A_k(A)^T\bar\Phi.
\]
\begin{theorem}\label{thm:fermion-block-integration}
Assume the hypotheses of \cref{thm:chamber-locality} and
\eqref{eq:weakened-walk-contraction}, with
$\kappa_{\mathrm w}=\kappa+\kappa_0+3$ and $\kappa,\kappa_0>0$.
Let $A\in\widetilde{\mathcal G}_k$ and $0<\mu\le1$, with fermions
antiperiodic in every coordinate. Then $K_k$ is invertible.
For every Grassmann polynomial $F$ and external coarse fields
$\Phi,\bar\Phi$, the weighted integral is
\begin{multline}\label{eq:Gaussian-block-integration}
 \int\mathrm D\Psi\,
 e^{-\langle\bar\Psi,\mathcal D_k\Psi\rangle
 -b_L\langle\bar\Phi-\mathcal A_k(-A)\bar\Psi,
                  \Phi-\mathcal A_k(A)\Psi\rangle_L}
 F(\bar\Psi,\Psi)\\
 =(\det K_k)e^{-\langle\bar\Phi,D_{k+1}^0\Phi\rangle_L}
 \int\dd\mu_I(W)
 F(\bar\Psi^c+\bar W,\Psi^c+\Gamma_kW).
\end{multline}
The incoming measure is one per unit-site spin component for $k\ge1$,
and $1/48$ per chamber spin component for $k=0$; the latter
integration includes all $192$ components per photon site.
The target pairing uses measure $L^3$. The identity-Gaussian
integration weights are respectively $1$ and $\sqrt{48}$.

There are analytic functions $E_k^d(X,A)$, depending only on bonds
meeting the polymer $X$, such that
\begin{equation}\label{eq:main-determinant-quotient}
 \frac{\det K_k(A)}{\det K_k(0)}
   =\exp\sum_{X\in\mathcal P_k}
           \bigl(E_k^d(X,A)-E_k^d(X,0)\bigr),\qquad
 |E_k^d(X,A)|\le C_dM^3e^{-(\kappa+1)d_M(X)}.
\end{equation}
They are invariant under real gauge transformations, even in $A$,
and satisfy
$\overline{E_k^d(X,\overline A)}=E_k^d(X,A)$.
They are covariant under coordinate sign changes and translations
in $LM\Z^3$. The constant $C_d$ depends on the fixed
$\rho,L,\alpha,M,\kappa,\kappa_0$, and is independent of the level,
volume, connection in the stated domain and $\mu\in(0,1]$.
\end{theorem}

The block identity follows by completing the quadratic form with the
weighted transposes. Its determinant is then localized by the
resolvent integral in \cref{subsec:determinant-integral}; the polymer
estimates permit that integral to be taken term by term.

For $\Psi=\Psi^c+W'$ and $\bar\Psi=\bar\Psi^c+\bar W'$,
completion of the square is the exact identity
\begin{multline}\label{eq:block-square-completion}
 \langle\bar\Psi,\mathcal D_k\Psi\rangle
   +b_L\langle\bar\Phi-\mathcal A_k(-A)\bar\Psi,
                    \Phi-\mathcal A_k(A)\Psi\rangle_L\\
   =\langle\bar\Phi,D_{k+1}^0\Phi\rangle_L
                    +\langle\bar W',K_kW'\rangle.
\end{multline}
Indeed, write $s=b_L\mathcal A_k^\sharp\Phi$ and
$\bar s=b_L\mathcal A_k(A)^T\bar\Phi$. The left side is
\[
 \langle\bar\Psi,K_k\Psi\rangle
      -\langle\bar\Psi,s\rangle-\langle\bar s,\Psi\rangle
      +b_L\langle\bar\Phi,\Phi\rangle_L.
\]
Substitute the means $\Gamma_ks,\Gamma_k^T\bar s$. Both cross
terms cancel, and the remaining constant is
$b_L\langle\bar\Phi,\Phi\rangle_L
-\langle\bar s,\Gamma_ks\rangle
=\langle\bar\Phi,D_{k+1}^0\Phi\rangle_L$.
All factors remain in their displayed order; no conjugation of
Grassmann fields is involved.

Translation by external odd fields preserves the top coefficient
in the integration variables. Applying this to
\eqref{eq:block-square-completion} and then using
\eqref{eq:covariance-substitution} proves
\eqref{eq:Gaussian-block-integration}, with $\delta Z_k(A)=\det K_k(A)$.
At the first step this is a single integration over all chamber
components. The determinant and the integration weight are those of
\eqref{eq:level-zero-covariance} and
\eqref{eq:first-step-Grassmann-weight}; the heavy components are
included in both.

\section{The local determinant}\label{sec:local-determinant}

\subsection{An integral for the determinant}\label{subsec:determinant-integral}

The determinant in \eqref{eq:Gaussian-block-integration} has a local
logarithm obtained by inserting the polymer propagators into a
resolvent integral. We first derive the integral at real connections.
Put $n_k=\ell^3$ and
\[
 c_k^d=4\bigl((1-L^{-3})\log b_k+L^{-3}\log(b_k+b_L)\bigr)
       \quad(k\ge1),\qquad c_0^d=0.
\]
For real $A\in\widetilde{\mathcal G}_k$ and $0<\mu\le1$,
\begin{align}
 \delta Z_k(A)
  &=\exp\left\{c_k^dn_k-ib_k^2\int_0^\infty\dd y\,
       \operatorname{Tr}\bigl[\gamma_5\mathcal B_{k,y}
          \mathcal Q_kS_y\mathcal Q_k^\sharp\mathcal B_{k,y}\bigr]
                        \right\}\quad(k\ge1),
                                    \label{eq:determinant-integral-positive-level}\\
 \delta Z_0(A)
  &=\exp\left\{\int_1^\infty\frac{\dd y}{y}
                       \operatorname{Tr}[F_0S_y]
             -i\int_0^1\dd y\,\operatorname{Tr}[\gamma_5S_y]
                        \right\}.\label{eq:determinant-integral-level-zero}
\end{align}
The trace in the first line is on the unit spinor space. At level
zero it is the normalized chamber trace
$\operatorname{Tr}O=48^{-1}\sum_{c,g}\tr O((c,g),(c,g))$.

To prove these identities, write $K=K_k$ as in the preceding
subsection. At real $A$, the fine adjoint identity
$D_A^*=\gamma_5D_A\gamma_5$, the reality of $\mu,b_k,b_L$, and
the orthogonal averaging projections give
$K^*=\gamma_5K\gamma_5$. Hence $K\gamma_5$ is Hermitian and
invertible. Since $\det\gamma_5=1$ on each four-component spinor,
its determinant is $\det K$. Let $\lambda_j$ be its real nonzero
eigenvalues, counted with multiplicity. There are $4n_k$ at positive
levels and $192n_0$ at level zero. Then
\[
 \operatorname{Tr}\bigl[\gamma_5(K+i\gamma_5y)^{-1}\bigr]
        =\sum_j(\lambda_j+iy)^{-1},\qquad
 \operatorname{Tr}\bigl[K(K+i\gamma_5y)^{-1}\bigr]
        =\sum_j\lambda_j(\lambda_j+iy)^{-1}.
\]
For a real $\lambda\ne0$, choose along $y\ge0$ the logarithm
with value $\log|\lambda|+i\pi\mathbf1[\lambda<0]$ at zero.
Its derivative is $i/(\lambda+iy)$ and
$\operatorname{Log}(\lambda+iy)-\log y\to i\pi/2$.
Integration gives
\begin{align*}
 -i\int_0^Y\frac{\dd y}{\lambda+iy}
        &=\operatorname{Log}\lambda-
                          \operatorname{Log}(\lambda+iY),\\
 \int_1^\infty\frac{\dd y}{y}\frac{\lambda}{\lambda+iy}
       -i\int_0^1\frac{\dd y}{\lambda+iy}
        &=\operatorname{Log}\lambda-i\pi/2.
\end{align*}
Summing the second identity proves
\eqref{eq:determinant-integral-level-zero}: its constant phase is
$-96\pi i n_0\in2\pi i\Z$.

For $k\ge1$, the ranks of $P_L$ and $1-P_L$ are respectively
$4L^{-3}n_k$ and $4(1-L^{-3})n_k$. The explicit formula for
$\mathcal B_{k,y}$ therefore gives
\[
 -i\int_0^Y\dd y\,\operatorname{Tr}[\gamma_5\mathcal B_{k,y}]
          =-4n_k\log Y+c_k^dn_k+o(1)
                  \quad(Y\to\infty).
\]
Insert \eqref{eq:coarse-resolvent-formula} in the first scalar
integral. The terms $4n_k\log Y$ cancel. The remaining phase is
$2\pi i n_k$, which disappears on exponentiation, and this proves
\eqref{eq:determinant-integral-positive-level}.
The integral converges absolutely: block compression bounds
$\mathcal Q_kS_y\mathcal Q_k^\sharp$ uniformly, while each row
and column sum of $\mathcal B_{k,y}$ is at most
$(1+2\exp(1))(b_k^2+y^2)^{-1/2}$. Its trace is bounded by
\[
 128\exp(1)(1+2\exp(1))^2a'n_k(b_k^2+y^2)^{-1}.
\]
At level zero the spectral identities already give absolute
convergence. The uniform kernel bounds below will give a bound
independent of the eigenvalues.

\subsection{Local determinant terms}\label{subsec:local-determinant-terms}

Choose $\gamma_{\mathrm w}\le\gamma$ and
$\kappa_{\mathrm w}=\kappa+\kappa_0+3$ with $\kappa,\kappa_0>0$.
The propagator bounds above give a fixed constant $C_w$ such that,
in component sup norms, for $y\ge0$,
\begin{align*}
 \|\mathbf1_{\Delta_z}S_{k,y}(X,A)\mathbf1_{\Delta_{z'}}\|
       &\le C_we^{-\gamma_{\mathrm w}d(z,z')
                                  -\kappa_{\mathrm w}d_M(X)}
                                                    \quad(k\ge1),\\
 |S_{0,y}(X,A)(x,x')|_c
       &\le C_w(1+y)^{-1}e^{-\gamma_{\mathrm w}d(c,c')
                                  -\kappa_{\mathrm w}d_M(X)}.
\end{align*}
Here $|\cdot|_c$ is the matrix norm induced by the component sup
norm. We may use the same $C_w$ after reducing the rates from
$\gamma,\kappa'$.
For $X\in\mathcal P_k$ define
\begin{align}
 \widetilde E_k(X,A)
    &=-ib_k^2\int_0^\infty\dd y\,
       \operatorname{Tr}\bigl[\gamma_5\mathcal B_{k,y}\mathcal Q_k
           S_{k,y}(X,A)\mathcal Q_k^\sharp\mathcal B_{k,y}\bigr]
                   \quad(k\ge1),\label{eq:local-determinant-definition}\\
 \widetilde E_0(X,A)
    &=\int_1^\infty\frac{\dd y}{y}
                     \operatorname{Tr}[F_0S_{0,y}(X,A)]
       -i\int_0^1\dd y\,\operatorname{Tr}[\gamma_5S_{0,y}(X,A)],\notag\\
 E_k^d(X,A)&=\tfrac14\bigl(\widetilde E_k(X,A)
       +\widetilde E_k(X,-A)
       +\overline{\widetilde E_k(X,\overline A)}
       +\overline{\widetilde E_k(X,-\overline A)}\bigr),\notag\\
 E_k^{\mathrm{det}}(X,A)&=E_k^d(X,A)-E_k^d(X,0).\notag
\end{align}
Set $c_X=54$ and
\[
 C_d=16c_XC_w\max\left\{\tfrac\pi2b\exp(1)(1+2\exp(1)),\,
                                  2+58\exp(1)+b\exp(1)\right\}.
\]

\begin{proposition}\label{prop:local-determinant-expansion}
The integrals \eqref{eq:local-determinant-definition} converge
absolutely for $A\in\widetilde{\mathcal G}_k$. They define analytic
functions, depending only on bonds meeting $X$, with
\begin{equation}\label{eq:local-determinant-bound}
 |\widetilde E_k(X,A)|, |E_k^d(X,A)|
               \le C_dM^3e^{-(\kappa+1)d_M(X)}.
\end{equation}
They satisfy
\begin{equation}\label{eq:local-determinant-quotient}
 \delta Z_k(A)=e^{c_k^dn_k+\sum_X\widetilde E_k(X,A)},\qquad
 \frac{\delta Z_k(A)}{\delta Z_k(0)}
                         =e^{\sum_XE_k^{\mathrm{det}}(X,A)}.
\end{equation}
The functions $E_k^d$ are gauge invariant, even in $A$, and satisfy
$\overline{E_k^d(X,\overline A)}=E_k^d(X,A)$. They are covariant
under the lattice step group.
\end{proposition}

\begin{proof}
For $k\ge1$ write
$N_y(X)=\mathcal Q_kS_{k,y}(X,A)\mathcal Q_k^\sharp$.
It is supported on the unit sites of $X$, and each entry is bounded
by $\exp(1)C_we^{-\kappa_{\mathrm w}d_M(X)}$ by normalized block
compression. To estimate the trace, cycle one factor
$\mathcal B_{k,y}$ to the left and use
\[
 G_y=\mathcal B_{k,y}\gamma_5\mathcal B_{k,y}
     =\gamma_5\bigl((b_k+i\gamma_5y)^{-2}(1-P_L)
              +(b_k+b_L+i\gamma_5y)^{-2}P_L\bigr).
\]
Its column sum in Euclidean matrix norm is at most
$(1+2\exp(1))/(b_k^2+y^2)$: the identity contributes one and the two
projection coefficients contribute at most $2\exp(1)$. Since
$|\tr(mm')|\le8|m|_2\max_{i,j}|m'_{ij}|$ for four-component
matrices, the integrand is bounded by
\begin{equation}\label{eq:polymer-determinant-trace-bound}
 8\exp(1)(1+2\exp(1))C_wM^3|X|_M
           e^{-\kappa_{\mathrm w}d_M(X)}(b_k^2+y^2)^{-1}.
\end{equation}
Integrating and multiplying by $b_k^2$ leaves the factor
$\pi b_k/2\le\pi b/2$.

At $k=0$, the formula for $D_A$ and the averaging penalty give
the scalar component row bound
\begin{equation}\label{eq:first-step-Dirac-row-bound}
 \sup_{x,\beta}\sum_{\beta'}\int\dd x'\,
                    |F_0(A)(x,x')_{\beta\beta'}|
                      \le R_F:=58\exp(1)+1+b\exp(1).
\end{equation}
The diagonal and three internal hops contribute
at most $41+\mu$, the external hop at most
$6e^{\mathfrak i(A)}$, and the block penalty at most $b\exp(1)$.
The bound follows from $\mu\le1$ and
$e^{\mathfrak i(A)}\le \exp(1)$. Point evaluation has normalized
kernel $48\mathbf1[x=x']$; its factor $48$ cancels the chamber
measure. The two traces in the level-zero integral are consequently
bounded by
\[
 4R_FC_wM^3|X|_M(1+y)^{-1}e^{-\kappa_{\mathrm w}d_M(X)},
 \qquad
 4C_wM^3|X|_M(1+y)^{-1}e^{-\kappa_{\mathrm w}d_M(X)}.
\]
The second bound uses that $\gamma_5$ has four entries of modulus
one. Their integrals contribute respectively
$\int_1^\infty(y(1+y))^{-1}\dd y=\log2$ and
$\int_0^1(1+y)^{-1}\dd y=\log2$.

The geometric bound \eqref{eq:polymer-size-tree} gives
$|X|_M\le c_X(d_M(X)+1)\le c_Xe^{d_M(X)}$.
Applying this to the trace bounds,
and using $\kappa_{\mathrm w}-1\ge\kappa+1$, proves
\eqref{eq:local-determinant-bound} for $\widetilde E_k$ and hence
for its four-term mean. The majorants are independent of $A$ and
integrable in $y$. Each integrand is analytic in $A$; Cauchy's
formula on a compact polydisc and the majorant permit its contour
integrals to be exchanged with the $y$ integral. The result is
analytic.

Every phase in the compressed trace belongs to a unit block or an
$L$-block meeting $X$, and that block is contained in $X$ by
nesting. The fine propagator is supported on $X$ and uses only its
connection. At level zero a nonzero external-hop trace uses a bond
meeting $X$; the penalty uses the $L$-block of its row. This proves
locality. Gauge transformations conjugate each product under the
trace and hence leave it unchanged. A lattice transformation whose
spin matrix commutes with $\gamma_5$ does the same. Every spatial
action of the step group has such a lift: multiplication by
$\gamma_1\gamma_2\gamma_3$ changes its commutation sign with
$\gamma_5$ and leaves its spatial action fixed. This proves step
covariance. The mean in \eqref{eq:local-determinant-definition}
then makes $E_k^d$ even and real in the stated analytic sense.

For real connections, sum the polymer propagators and use
\eqref{eq:determinant-integral-positive-level} or
\eqref{eq:determinant-integral-level-zero}. This proves the first
identity of \eqref{eq:local-determinant-quotient} there. Both sides
are analytic on $\widetilde{\mathcal G}_k$, which is open and
star-shaped. They agree on a real neighborhood of zero, so their
Taylor series agree there and analytic continuation proves the
identity throughout the domain.

To justify the symmetrized quotient, put
$\mathcal L(A)=\sum_X\widetilde E_k(X,A)$.
The charge-conjugation transpose and the complex-conjugation
identity give
$\delta Z_k(-A)=\delta Z_k(A)
=\overline{\delta Z_k(\overline A)}$.
For the latter, the spin matrix $\varkappa=\gamma_5\mathsf C$
conjugates $\overline{K_k(\overline A)}$ to $K_k(-A)$.
Thus $\mathcal L(-A)-\mathcal L(A)$ and
$\overline{\mathcal L(\overline A)}-\mathcal L(A)$ take values
in $2\pi i\Z$. They are continuous on the connected domain,
so they are constant. Subtracting their values at zero gives
\[
 \sum_X\bigl(E_k^d(X,A)-E_k^d(X,0)\bigr)
                            =\mathcal L(A)-\mathcal L(0).
\]
This proves the second identity in
\eqref{eq:local-determinant-quotient} without a choice of logarithm
branch for the complex determinant.
\end{proof}

\subsection{The small-field gain for the determinant}\label{subsec:determinant-small-field-gain}

For a photon connection write $A_\mu(c)=A(c,c+\eta_ke_\mu)$ and
\[
 \partial_\nu A_\mu(c)
       =\eta_k^{-1}(A_\mu(c+\eta_ke_\nu)-A_\mu(c)),\qquad
 \|\delta_\alpha\partial A\|_\infty
       =\sup_{\mu,\nu}\sup_{0<d(c,c')\le1}
          \frac{|\partial_\nu A_\mu(c')-\partial_\nu A_\mu(c)|}
               {d(c,c')^\alpha}.
\]
Local versions use only differences and bonds inside the given set.
In this section the curl is mesh normalized:
$(dA)_{\mu\nu}=\partial_\mu A_\nu-\partial_\nu A_\mu$.
Let $R_k$ be the complex bond functions whose three norms
$\|A\|_\infty$, $\|\partial A\|_\infty$, and
$\|\delta_\alpha\partial A\|_\infty$ are less than
$e_k^{-3/4+\varepsilon}$, $e_k^{-3/4+2\varepsilon}$, and
$e_k^{-3/4+3\varepsilon}$, respectively. Let $\widetilde R_k$
consist of $A_0+A_1$, with $A_1\in R_k$ and $A_0$ real, such that
on every $\square^\dagger=\square^{\sim L}$ one can write
$A_0=A'+d\lambda$, where $\lambda,A'$ are real and the same
three local bounds hold for $A'$. The external connection domain is
\[
 \mathfrak X_k=\{\mathscr C:\|\mathscr C\|_\infty
                  <e_k^{-1+2\varepsilon},\quad
             \|\partial\mathscr C\|_\infty<e_k^{-1+3\varepsilon}\}.
\]
If $e_k^{1/4}+e_k^\varepsilon<1$, both $\mathfrak X_k$ and
$\widetilde R_k+\mathfrak X_k$ lie in
$\widetilde{\mathcal G}_k$: for $\mathscr A=A_0+A_1\in\widetilde R_k$,
the real summand $A_0$ has on the local cubes a representative
of size $e_k^{-3/4+\varepsilon}$, and the remainder
$A_1+\mathscr C$ has size at most
$(e_k^{1/4}+e_k^\varepsilon)e_k^{-1+\varepsilon}$.
For $\mathfrak X_k$ alone take $A_0=A_1=0$.

Put $v_k=-\log e_k$. Assume
\begin{equation}\label{eq:determinant-gain-conditions}
 v_k\ge9,\qquad e_k^\varepsilon\le\tfrac14,\qquad
 8M(v_k/4+4)e_k^\varepsilon\le1,\qquad
 \ell\ge Mv_k^2.
\end{equation}

\begin{proposition}\label{prop:determinant-small-field-gain}
For $\mathscr A\in\widetilde R_k$, $\mathscr C\in\mathfrak X_k$,
and $X\in\mathcal P_k$, assume \eqref{eq:determinant-gain-conditions}
and $e_k^{1/4}+e_k^\varepsilon<1$. Then
\begin{equation}\label{eq:determinant-small-field-gain}
 |E_k^d(X,\mathscr A+\mathscr C)-E_k^d(X,\mathscr C)|
       \le8C_dM^3e_k^{1/4}e^{-\kappa d_M(X)}.
\end{equation}
\end{proposition}

\begin{proof}
If $d_M(X)>v_k/4$, bound the two terms separately by
\eqref{eq:local-determinant-bound} and use
$e^{-d_M(X)}\le e_k^{1/4}$. Suppose therefore that
$d_M(X)\le v_k/4$. The polymer lift and diameter estimates give
a nonwrapping lattice box $W$ containing every bond meeting $X$,
with each side of length at most
$s_W=M(d_M(X)+3)$.
Their applicability follows from
$2M(d_M(X)+2)<\ell$ and $M(d_M(X)+4)<\ell$, which hold by
\eqref{eq:determinant-gain-conditions}.

Write $\mathscr A=A_0+A_1$ as in the domain definition.
On each plaquette, the local real representative gives
$\|dA_0\|_\infty<2e_k^{-3/4+2\varepsilon}$.
Choose a fixed-order rectilinear path $\Gamma(o,c)$ from a site
$o\in W$ and set $\lambda(c)=\eta_k A_0(\Gamma(o,c))$ on $W$.
Then
\[
 \|A_0-d\lambda\|_{\infty,W}
             \le2s_W\|dA_0\|_{\infty,W}
             <4s_We_k^{-3/4+2\varepsilon}.
\]
To see the factor $2s_W$, append one bond to $\Gamma(o,c)$ and
subtract $\Gamma(o,c')$. After cancelling common bonds, the
resulting loop bounds at most two rectangles of width $\eta_k$.
Together they contain at most $2s_W/\eta_k$ plaquettes, each
contributing $\eta_k dA_0$. The paths and rectangles remain in
$W$ because it is a product of intervals. Extend the real function
$\lambda$ by zero outside the projected box. On every bond meeting
$X$, the bound is unchanged. By
\eqref{eq:determinant-gain-conditions} it is at most
$\tfrac12e_k^{-3/4+\varepsilon}$.

Let $B=A_0-d\lambda+A_1$ on bonds meeting $X$ and zero elsewhere.
Then $\|B\|_\infty<\tfrac32e_k^{-3/4+\varepsilon}$.
For $|t|\le\tfrac13e_k^{-1/4}$,
\[
 \|tB+\mathscr C\|_\infty
       <(\tfrac12+e_k^\varepsilon)e_k^{-1+\varepsilon}
       \le\tfrac34e_k^{-1+\varepsilon}.
\]
Thus this connection belongs to the complex connection ball
$\mathcal G_k=\{A:\|A\|_\infty<e_k^{-1+\varepsilon}\}$,
which is contained in $\widetilde{\mathcal G}_k$, and
$f(t)=E_k^d(X,tB+\mathscr C)$ is analytic on this disk and
bounded by the right side of \eqref{eq:local-determinant-bound}.
Gauge invariance and locality identify the difference in
\eqref{eq:determinant-small-field-gain} with $f(1)-f(0)$.
On the circle of radius $r=\tfrac14e_k^{-1/4}>2$, Cauchy's
formula yields
\[
 f(1)-f(0)=\frac1{2\pi i}\oint_{|t|=r}
                         \frac{f(t)}{t(t-1)}\,\dd t,
 \qquad |f(1)-f(0)|\le\frac{\sup|f|}{r-1}
                          \le8e_k^{1/4}\sup|f|.
\]
This proves the claimed gain, retaining the stronger decay
$e^{-(\kappa+1)d_M(X)}$ in this case.
\end{proof}

\begin{proof}[Proof of \Cref{thm:fermion-block-integration}]
\Cref{prop:fine-coarse-resolvents} and
\eqref{eq:level-zero-covariance} give invertibility on the two
integration spaces. The weighted completion calculation gives
\eqref{eq:Gaussian-block-integration}, with the first-step weight from
\eqref{eq:first-step-Grassmann-weight}.
\Cref{prop:local-determinant-expansion} gives the analytic determinant
terms, their bound, symmetries and normalized quotient.
\end{proof}


\appendix
\section{Polymer geometry}\label{app:polymer-geometry}

We use the geometric arguments of \cite[Appendix~A]{QED3Manuscript}.
Let $L$ be odd, $M$ a power of $L$, and $\ell$ a multiple of $LM$
with $\ell\ge13LM$. The $M$-cubes have centers in
$M\Z^3/\ell\Z^3$. They partition the lattice sites; their closed
continuum versions may meet along faces and edges. A polymer is a
nonempty union of $M$-cubes connected through faces. Denote its number
of cubes by $|X|_M$. A net in the continuum torus is a connected
finite union of images of specified closed straight segments
in $\R^3$, a point being allowed as a segment.
Its length is the sum of the Euclidean lengths of those segments.
Set
\[
 Md_M(X)=\inf\{\operatorname{length}(N):N
             \text{ is a net meeting every closed cube of }X\},
\]
where the infimum also extends over the finite segment representations.
The same definition is used for finite face-connected unions of
$M$-cubes in $\R^3$, with nets in $\R^3$. Distances are sup distances.


\subsection{Size and diameter}
The following bounds hold on the torus and on the covering space:
\begin{align}
 d_M(X)&\le |X|_M-1,&
 |X|_M&\le54(d_M(X)+1),\label{eq:polymer-size-tree}\\
 d(p,q)&\le M(d_M(X)+2)
       &&(p,q\text{ in the closed cubes of }X).
                                  \label{eq:polymer-diameter}
\end{align}

To prove the first inequality of \eqref{eq:polymer-size-tree}, order
the cubes along a spanning tree of their adjacency graph. A net
meeting one cube can be extended to an adjacent cube by a segment of
length at most $M$ perpendicular to their common face. Starting with
one point and repeating this construction gives a net of length at
most $M(|X|_M-1)$.

For the reverse inequality, subdivide the segments of any net $N$ at
their intersections and discard overlapping pieces. Its total length
does not increase. Doubling the edges of the resulting connected
graph gives a closed traversal of length at most
$2\operatorname{length}(N)$. Cut it into at most
$2\operatorname{length}(N)/M+1$ arcs of length at most $M$.
Each arc lifts to a curve in $\R^3$ whose range in each coordinate
has length at most $M$. Such an interval meets at most three closed
cube intervals, so an arc meets at most $27$ closed $M$-cubes.
Every cube of $X$ meets one of these arcs. Taking the infimum over
$N$ gives
\[
 |X|_M\le27(2d_M(X)+1)\le c_X(d_M(X)+1).
\]
For \eqref{eq:polymer-diameter}, take a net of length less than
$Md_M(X)+\delta$ and points $p',q'$ of the net in the respective
cubes of $p,q$. The path in the net from $p'$ to $q'$ has length
at most that of the net. Hence
$d(p,q)\le M+(Md_M(X)+\delta)+M$; let $\delta\downarrow0$.
The same argument holds in $\R^3$.

\subsection{Nonwrapping lifts}\label{subsec:polymer-lifts}

Let $\operatorname{pr}:\R^3\to\R^3/\ell\Z^3$ be the projection.
\begin{proposition}\label{prop:nonwrapping-polymer-lift}
If $X$ is a polymer and
\begin{equation}\label{eq:nonwrapping-polymer-condition}
 2M(d_M(X)+2)<\ell,
\end{equation}
then $X=\operatorname{pr}(\widetilde X)$ for a finite face-connected
union of $M$-cubes in $\R^3$. The projection is injective on
$\widetilde X$, preserves adjacency of its cubes in both directions,
and $d_M(\widetilde X)=d_M(X)$. The lift is unique up to a translation
in $\ell\Z^3$.
\end{proposition}

\begin{proof}
Put $D=M(d_M(X)+2)$, fix $p_0$ in a closed cube of $X$, and fix its
lift $\widetilde p_0$. The projection is injective on the closed sup
ball $\widetilde B=\{|p-\widetilde p_0|_\infty\le D\}$, since two
points there differ in every coordinate by at most $2D<\ell$.
It maps this ball onto the torus ball $B=\{p:d(p,p_0)\le D\}$.
Its inverse $\varphi:B\to\widetilde B$ is continuous, and is locally
isometric. The diameter estimate places every closed cube of $X$ in
$B$. A segment or a closed $M$-cube lying in $B$ lifts under
$\varphi$ to a segment of the same length or a closed $M$-cube:
the translates $\widetilde B+\ell n$ are pairwise disjoint, and a
connected lift lies in one of them. The same inverse maps shared
faces to shared faces. Thus $\widetilde X=\varphi(X)$ has the
asserted injectivity, adjacency, and connectedness.

Any other face-connected injective lift can be translated so that it
contains the chosen lift of the cube of $p_0$. There are no two
different neighbors of one lifted cube with the same projection:
their centers differ by at most $2M<\ell$. Propagate this observation
along a chain of adjacent cubes. Every cube of the other lift equals
the corresponding cube of $\widetilde X$, proving uniqueness.

Projecting a net for $\widetilde X$ gives a net for $X$ with no greater
length, so $d_M(X)\le d_M(\widetilde X)$. Conversely, let $N$ be a
net for $X$ of length at most $Md_M(X)+M$, and choose its point $q_0$
in the cube of $p_0$. Every point of $N$ is joined to $q_0$ by a
curve of length at most that of $N$, and $d(q_0,p_0)\le M$.
Therefore $N\subset B$. Its lift $\varphi(N)$ has the same length
and meets every cube of $\widetilde X$. Taking an infimizing sequence
of these nets proves the reverse inequality.
\end{proof}

For the determinant estimate, suppose in addition that
$M(d_M(X)+4)<\ell$ and $M\ge2$. The closed cubes of $\widetilde X$
have coordinate ranges of length at most $M(d_M(X)+2)$ by
\eqref{eq:polymer-diameter}. Enlarge these intervals by one mesh
$\eta_k\le1$ at each end and intersect with $\eta_k\Z^3$. The
resulting box $W$ contains every bond meeting the lifted polymer and
has each side of length at most $M(d_M(X)+3)$. This length, increased
by one mesh, is less than $\ell$. Thus the projection is injective
on $W$, and every torus bond meeting $X$ has exactly one lift in $W$.

\printbibliography
\end{document}
